% concatenated from preprint.tex, math_commands.tex
\documentclass{article}

% ready for submission
\usepackage[preprint]{neurips_2026}
\usepackage{enumitem}
\usepackage[utf8]{inputenc}
\usepackage[T1]{fontenc}
\usepackage{hyperref}
\usepackage{url}
\usepackage{booktabs}
\usepackage{amsfonts}
\usepackage{nicefrac}
%\usepackage{microtype}
\usepackage{xcolor}
\usepackage{amsmath, amssymb, amsthm}
\usepackage{bm}
\usepackage{algorithm}
\usepackage{algorithmic}
\usepackage{graphicx}
%%%%% NEW MATH DEFINITIONS %%%%%



\usepackage{amsmath,amsfonts,bm}

% Mark sections of captions for referring to divisions of figures
\newcommand{\figleft}{{\em (Left)}}
\newcommand{\figcenter}{{\em (Center)}}
\newcommand{\figright}{{\em (Right)}}
\newcommand{\figtop}{{\em (Top)}}
\newcommand{\figbottom}{{\em (Bottom)}}
\newcommand{\captiona}{{\em (a)}}
\newcommand{\captionb}{{\em (b)}}
\newcommand{\captionc}{{\em (c)}}
\newcommand{\captiond}{{\em (d)}}

% Highlight a newly defined term
\newcommand{\newterm}[1]{{\bf #1}}


% Figure reference, lower-case.
\def\figref#1{figure~\ref{#1}}
% Figure reference, capital. For start of sentence
\def\Figref#1{Figure~\ref{#1}}
\def\twofigref#1#2{figures \ref{#1} and \ref{#2}}
\def\quadfigref#1#2#3#4{figures \ref{#1}, \ref{#2}, \ref{#3} and \ref{#4}}
% Section reference, lower-case.
\def\secref#1{section~\ref{#1}}
% Section reference, capital.
\def\Secref#1{Section~\ref{#1}}
% Reference to two sections.
\def\twosecrefs#1#2{sections \ref{#1} and \ref{#2}}
% Reference to three sections.
\def\secrefs#1#2#3{sections \ref{#1}, \ref{#2} and \ref{#3}}
% Reference to an equation, lower-case.
\def\eqref#1{equation~\ref{#1}}
% Reference to an equation, upper case
\def\Eqref#1{Equation~\ref{#1}}
% A raw reference to an equation---avoid using if possible
\def\plaineqref#1{\ref{#1}}
% Reference to a chapter, lower-case.
\def\chapref#1{chapter~\ref{#1}}
% Reference to an equation, upper case.
\def\Chapref#1{Chapter~\ref{#1}}
% Reference to a range of chapters
\def\rangechapref#1#2{chapters\ref{#1}--\ref{#2}}
% Reference to an algorithm, lower-case.
\def\algref#1{algorithm~\ref{#1}}
% Reference to an algorithm, upper case.
\def\Algref#1{Algorithm~\ref{#1}}
\def\twoalgref#1#2{algorithms \ref{#1} and \ref{#2}}
\def\Twoalgref#1#2{Algorithms \ref{#1} and \ref{#2}}
% Reference to a part, lower case
\def\partref#1{part~\ref{#1}}
% Reference to a part, upper case
\def\Partref#1{Part~\ref{#1}}
\def\twopartref#1#2{parts \ref{#1} and \ref{#2}}

\def\ceil#1{\lceil #1 \rceil}
\def\floor#1{\lfloor #1 \rfloor}
\def\1{\bm{1}}
\newcommand{\train}{\mathcal{D}}
\newcommand{\valid}{\mathcal{D_{\mathrm{valid}}}}
\newcommand{\test}{\mathcal{D_{\mathrm{test}}}}

\def\eps{{\epsilon}}
\newcommand{\cX}{\mathcal{X}}
\newcommand{\cO}{\mathcal{O}}

\def\beq{\begin{equation}}
\def\eeq{\end{equation}}
\def\ba{\begin{array}}
\def\ea{\end{array}}
\def\beann{\begin{eqnarray*}}
\def\eeann{\end{eqnarray*}}
\def\bea{\begin{eqnarray}}
\def\eea{\end{eqnarray}}

\def\BT{\begin{theorem}}
\def\ET{\end{theorem}}
\def\BL{\begin{lemma}}
\def\EL{\end{lemma}}
\def\BC{\begin{corollary}}
\def\EC{\end{corollary}}
\def\BE{\begin{example}}
\def\EE{\end{example}}
\def\BD{\begin{definition}}
\def\ED{\end{definition}}
\def\BR{\begin{remark}}
\def\ER{\end{remark}}
\def\BAS{\begin{assumption}}
\def\EAS{\end{assumption}}
\def\BI{\begin{itemize}}
\def\EI{\end{itemize}}
\def\BP{\begin{proposition}}
\def\EP{\end{proposition}}

\def\BMP{\begin{minipage}{9.5cm}}
\def\EMP{\end{minipage}}
\def\MPT{\begin{minipage}{11.5cm}}
\def\EPT{\end{minipage}}

\def\la{\langle}
\def\ra{\rangle}

\def\I{\mathbb{I}}

\def\QF{\hspace{5ex} \Box}
\def\QR{\hfill \Box}

% Random variables
\def\reta{{\textnormal{$\eta$}}}
\def\ra{{\textnormal{a}}}
\def\rb{{\textnormal{b}}}
\def\rc{{\textnormal{c}}}
\def\rd{{\textnormal{d}}}
\def\re{{\textnormal{e}}}
\def\rf{{\textnormal{f}}}
\def\rg{{\textnormal{g}}}
\def\rh{{\textnormal{h}}}
\def\ri{{\textnormal{i}}}
\def\rj{{\textnormal{j}}}
\def\rk{{\textnormal{k}}}
\def\rl{{\textnormal{l}}}
% rm is already a command, just don't name any random variables m
\def\rn{{\textnormal{n}}}
\def\ro{{\textnormal{o}}}
\def\rp{{\textnormal{p}}}
\def\rq{{\textnormal{q}}}
\def\rr{{\textnormal{r}}}
\def\rs{{\textnormal{s}}}
\def\rt{{\textnormal{t}}}
\def\ru{{\textnormal{u}}}
\def\rv{{\textnormal{v}}}
\def\rw{{\textnormal{w}}}
\def\rx{{\textnormal{x}}}
\def\ry{{\textnormal{y}}}
\def\rz{{\textnormal{z}}}

% Random vectors
\def\rvepsilon{{\mathbf{\epsilon}}}
\def\rvtheta{{\mathbf{\theta}}}
\def\rva{{\mathbf{a}}}
\def\rvb{{\mathbf{b}}}
\def\rvc{{\mathbf{c}}}
\def\rvd{{\mathbf{d}}}
\def\rve{{\mathbf{e}}}
\def\rvf{{\mathbf{f}}}
\def\rvg{{\mathbf{g}}}
\def\rvh{{\mathbf{h}}}
\def\rvu{{\mathbf{i}}}
\def\rvj{{\mathbf{j}}}
\def\rvk{{\mathbf{k}}}
\def\rvl{{\mathbf{l}}}
\def\rvm{{\mathbf{m}}}
\def\rvn{{\mathbf{n}}}
\def\rvo{{\mathbf{o}}}
\def\rvp{{\mathbf{p}}}
\def\rvq{{\mathbf{q}}}
\def\rvr{{\mathbf{r}}}
\def\rvs{{\mathbf{s}}}
\def\rvt{{\mathbf{t}}}
\def\rvu{{\mathbf{u}}}
\def\rvv{{\mathbf{v}}}
\def\rvw{{\mathbf{w}}}
\def\rvx{{\mathbf{x}}}
\def\rvy{{\mathbf{y}}}
\def\rvz{{\mathbf{z}}}

% Elements of random vectors
\def\erva{{\textnormal{a}}}
\def\ervb{{\textnormal{b}}}
\def\ervc{{\textnormal{c}}}
\def\ervd{{\textnormal{d}}}
\def\erve{{\textnormal{e}}}
\def\ervf{{\textnormal{f}}}
\def\ervg{{\textnormal{g}}}
\def\ervh{{\textnormal{h}}}
\def\ervi{{\textnormal{i}}}
\def\ervj{{\textnormal{j}}}
\def\ervk{{\textnormal{k}}}
\def\ervl{{\textnormal{l}}}
\def\ervm{{\textnormal{m}}}
\def\ervn{{\textnormal{n}}}
\def\ervo{{\textnormal{o}}}
\def\ervp{{\textnormal{p}}}
\def\ervq{{\textnormal{q}}}
\def\ervr{{\textnormal{r}}}
\def\ervs{{\textnormal{s}}}
\def\ervt{{\textnormal{t}}}
\def\ervu{{\textnormal{u}}}
\def\ervv{{\textnormal{v}}}
\def\ervw{{\textnormal{w}}}
\def\ervx{{\textnormal{x}}}
\def\ervy{{\textnormal{y}}}
\def\ervz{{\textnormal{z}}}

% Random matrices
\def\rmA{{\mathbf{A}}}
\def\rmB{{\mathbf{B}}}
\def\rmC{{\mathbf{C}}}
\def\rmD{{\mathbf{D}}}
\def\rmE{{\mathbf{E}}}
\def\rmF{{\mathbf{F}}}
\def\rmG{{\mathbf{G}}}
\def\rmH{{\mathbf{H}}}
\def\rmI{{\mathbf{I}}}
\def\rmJ{{\mathbf{J}}}
\def\rmK{{\mathbf{K}}}
\def\rmL{{\mathbf{L}}}
\def\rmM{{\mathbf{M}}}
\def\rmN{{\mathbf{N}}}
\def\rmO{{\mathbf{O}}}
\def\rmP{{\mathbf{P}}}
\def\rmQ{{\mathbf{Q}}}
\def\rmR{{\mathbf{R}}}
\def\rmS{{\mathbf{S}}}
\def\rmT{{\mathbf{T}}}
\def\rmU{{\mathbf{U}}}
\def\rmV{{\mathbf{V}}}
\def\rmW{{\mathbf{W}}}
\def\rmX{{\mathbf{X}}}
\def\rmY{{\mathbf{Y}}}
\def\rmZ{{\mathbf{Z}}}

% Elements of random matrices
\def\ermA{{\textnormal{A}}}
\def\ermB{{\textnormal{B}}}
\def\ermC{{\textnormal{C}}}
\def\ermD{{\textnormal{D}}}
\def\ermE{{\textnormal{E}}}
\def\ermF{{\textnormal{F}}}
\def\ermG{{\textnormal{G}}}
\def\ermH{{\textnormal{H}}}
\def\ermI{{\textnormal{I}}}
\def\ermJ{{\textnormal{J}}}
\def\ermK{{\textnormal{K}}}
\def\ermL{{\textnormal{L}}}
\def\ermM{{\textnormal{M}}}
\def\ermN{{\textnormal{N}}}
\def\ermO{{\textnormal{O}}}
\def\ermP{{\textnormal{P}}}
\def\ermQ{{\textnormal{Q}}}
\def\ermR{{\textnormal{R}}}
\def\ermS{{\textnormal{S}}}
\def\ermT{{\textnormal{T}}}
\def\ermU{{\textnormal{U}}}
\def\ermV{{\textnormal{V}}}
\def\ermW{{\textnormal{W}}}
\def\ermX{{\textnormal{X}}}
\def\ermY{{\textnormal{Y}}}
\def\ermZ{{\textnormal{Z}}}

% Vectors
\def\vzero{{\bm{0}}}
\def\vone{{\bm{1}}}
\def\vmu{{\bm{\mu}}}
\def\vtheta{{\bm{\theta}}}
\def\va{{\bm{a}}}
\def\vb{{\bm{b}}}
\def\vc{{\bm{c}}}
\def\vd{{\bm{d}}}
\def\ve{{\bm{e}}}
\def\vf{{\bm{f}}}
\def\vg{{\bm{g}}}
\def\vh{{\bm{h}}}
\def\vi{{\bm{i}}}
\def\vj{{\bm{j}}}
\def\vk{{\bm{k}}}
\def\vl{{\bm{l}}}
\def\vm{{\bm{m}}}
\def\vn{{\bm{n}}}
\def\vo{{\bm{o}}}
\def\vp{{\bm{p}}}
\def\vq{{\bm{q}}}
\def\vr{{\bm{r}}}
\def\vs{{\bm{s}}}
\def\vt{{\bm{t}}}
\def\vu{{\bm{u}}}
\def\vv{{\bm{v}}}
\def\vw{{\bm{w}}}
\def\vx{{\bm{x}}}
\def\vy{{\bm{y}}}
\def\vz{{\bm{z}}}
\def\vX{{\bm{X}}}
\def\vR{{\bm{R}}}
\def\vA{{\bm{A}}}
\def\vB{{\bm{B}}}
\def\vV{{\bm{V}}}
% Elements of vectors
\def\evalpha{{\alpha}}
\def\evbeta{{\beta}}
\def\evepsilon{{\epsilon}}
\def\evlambda{{\lambda}}
\def\evomega{{\omega}}
\def\evmu{{\mu}}
\def\evpsi{{\psi}}
\def\evsigma{{\sigma}}
\def\evtheta{{\theta}}
\def\eva{{a}}
\def\evb{{b}}
\def\evc{{c}}
\def\evd{{d}}
\def\eve{{e}}
\def\evf{{f}}
\def\evg{{g}}
\def\evh{{h}}
\def\evi{{i}}
\def\evj{{j}}
\def\evk{{k}}
\def\evl{{l}}
\def\evm{{m}}
\def\evn{{n}}
\def\evo{{o}}
\def\evp{{p}}
\def\evq{{q}}
\def\evr{{r}}
\def\evs{{s}}
\def\evt{{t}}
\def\evu{{u}}
\def\evv{{v}}
\def\evw{{w}}
\def\evx{{x}}
\def\evy{{y}}
\def\evz{{z}}

% Matrix
\def\mA{{\bm{A}}}
\def\mB{{\bm{B}}}
\def\mC{{\bm{C}}}
\def\mD{{\bm{D}}}
\def\mE{{\bm{E}}}
\def\mF{{\bm{F}}}
\def\mG{{\bm{G}}}
\def\mH{{\bm{H}}}
\def\mI{{\bm{I}}}
\def\mJ{{\bm{J}}}
\def\mK{{\bm{K}}}
\def\mL{{\bm{L}}}
\def\mM{{\bm{M}}}
\def\mN{{\bm{N}}}
\def\mO{{\bm{O}}}
\def\mP{{\bm{P}}}
\def\mQ{{\bm{Q}}}
\def\mR{{\bm{R}}}
\def\mS{{\bm{S}}}
\def\mT{{\bm{T}}}
\def\mU{{\bm{U}}}
\def\mV{{\bm{V}}}
\def\mW{{\bm{W}}}
\def\mX{{\bm{X}}}
\def\mY{{\bm{Y}}}
\def\mZ{{\bm{Z}}}
\def\mBeta{{\bm{\beta}}}
\def\mPhi{{\bm{\Phi}}}
\def\mLambda{{\bm{\Lambda}}}
\def\mSigma{{\bm{\Sigma}}}

% Tensor
\DeclareMathAlphabet{\mathsfit}{\encodingdefault}{\sfdefault}{m}{sl}
\SetMathAlphabet{\mathsfit}{bold}{\encodingdefault}{\sfdefault}{bx}{n}
\newcommand{\tens}[1]{\bm{\mathsfit{#1}}}
\def\tA{{\tens{A}}}
\def\tB{{\tens{B}}}
\def\tC{{\tens{C}}}
\def\tD{{\tens{D}}}
\def\tE{{\tens{E}}}
\def\tF{{\tens{F}}}
\def\tG{{\tens{G}}}
\def\tH{{\tens{H}}}
\def\tI{{\tens{I}}}
\def\tJ{{\tens{J}}}
\def\tK{{\tens{K}}}
\def\tL{{\tens{L}}}
\def\tM{{\tens{M}}}
\def\tN{{\tens{N}}}
\def\tO{{\tens{O}}}
\def\tP{{\tens{P}}}
\def\tQ{{\tens{Q}}}
\def\tR{{\tens{R}}}
\def\tS{{\tens{S}}}
\def\tT{{\tens{T}}}
\def\tU{{\tens{U}}}
\def\tV{{\tens{V}}}
\def\tW{{\tens{W}}}
\def\tX{{\tens{X}}}
\def\tY{{\tens{Y}}}
\def\tZ{{\tens{Z}}}


% Graph
\def\gA{{\mathcal{A}}}
\def\gB{{\mathcal{B}}}
\def\gC{{\mathcal{C}}}
\def\gD{{\mathcal{D}}}
\def\gE{{\mathcal{E}}}
\def\gF{{\mathcal{F}}}
\def\gG{{\mathcal{G}}}
\def\gH{{\mathcal{H}}}
\def\gI{{\mathcal{I}}}
\def\gJ{{\mathcal{J}}}
\def\gK{{\mathcal{K}}}
\def\gL{{\mathcal{L}}}
\def\gM{{\mathcal{M}}}
\def\gN{{\mathcal{N}}}
\def\gO{{\mathcal{O}}}
\def\gP{{\mathcal{P}}}
\def\gQ{{\mathcal{Q}}}
\def\gR{{\mathcal{R}}}
\def\gS{{\mathcal{S}}}
\def\gT{{\mathcal{T}}}
\def\gU{{\mathcal{U}}}
\def\gV{{\mathcal{V}}}
\def\gW{{\mathcal{W}}}
\def\gX{{\mathcal{X}}}
\def\gY{{\mathcal{Y}}}
\def\gZ{{\mathcal{Z}}}

% Sets
\def\sA{{\mathbb{A}}}
\def\sB{{\mathbb{B}}}
\def\sC{{\mathbb{C}}}
\def\sD{{\mathbb{D}}}
% Don't use a set called E, because this would be the same as our symbol
% for expectation.
\def\sF{{\mathbb{F}}}
\def\sG{{\mathbb{G}}}
\def\sH{{\mathbb{H}}}
\def\sI{{\mathbb{I}}}
\def\sJ{{\mathbb{J}}}
\def\sK{{\mathbb{K}}}
\def\sL{{\mathbb{L}}}
\def\sM{{\mathbb{M}}}
\def\sN{{\mathbb{N}}}
\def\sO{{\mathbb{O}}}
\def\sP{{\mathbb{P}}}
\def\sQ{{\mathbb{Q}}}
\def\sR{{\mathbb{R}}}
\def\sS{{\mathbb{S}}}
\def\sT{{\mathbb{T}}}
\def\sU{{\mathbb{U}}}
\def\sV{{\mathbb{V}}}
\def\sW{{\mathbb{W}}}
\def\sX{{\mathbb{X}}}
\def\sY{{\mathbb{Y}}}
\def\sZ{{\mathbb{Z}}}

% Entries of a matrix
\def\emLambda{{\Lambda}}
\def\emA{{A}}
\def\emB{{B}}
\def\emC{{C}}
\def\emD{{D}}
\def\emE{{E}}
\def\emF{{F}}
\def\emG{{G}}
\def\emH{{H}}
\def\emI{{I}}
\def\emJ{{J}}
\def\emK{{K}}
\def\emL{{L}}
\def\emM{{M}}
\def\emN{{N}}
\def\emO{{O}}
\def\emP{{P}}
\def\emQ{{Q}}
\def\emR{{R}}
\def\emS{{S}}
\def\emT{{T}}
\def\emU{{U}}
\def\emV{{V}}
\def\emW{{W}}
\def\emX{{X}}
\def\emY{{Y}}
\def\emZ{{Z}}
\def\emSigma{{\Sigma}}

% entries of a tensor
% Same font as tensor, without \bm wrapper
\newcommand{\etens}[1]{\mathsfit{#1}}
\def\etLambda{{\etens{\Lambda}}}
\def\etA{{\etens{A}}}
\def\etB{{\etens{B}}}
\def\etC{{\etens{C}}}
\def\etD{{\etens{D}}}
\def\etE{{\etens{E}}}
\def\etF{{\etens{F}}}
\def\etG{{\etens{G}}}
\def\etH{{\etens{H}}}
\def\etI{{\etens{I}}}
\def\etJ{{\etens{J}}}
\def\etK{{\etens{K}}}
\def\etL{{\etens{L}}}
\def\etM{{\etens{M}}}
\def\etN{{\etens{N}}}
\def\etO{{\etens{O}}}
\def\etP{{\etens{P}}}
\def\etQ{{\etens{Q}}}
\def\etR{{\etens{R}}}
\def\etS{{\etens{S}}}
\def\etT{{\etens{T}}}
\def\etU{{\etens{U}}}
\def\etV{{\etens{V}}}
\def\etW{{\etens{W}}}
\def\etX{{\etens{X}}}
\def\etY{{\etens{Y}}}
\def\etZ{{\etens{Z}}}

% The true underlying data generating distribution
\newcommand{\pdata}{p_{\rm{data}}}
% The empirical distribution defined by the training set
\newcommand{\ptrain}{\hat{p}_{\rm{data}}}
\newcommand{\Ptrain}{\hat{P}_{\rm{data}}}
% The model distribution
\newcommand{\pmodel}{p_{\rm{model}}}
\newcommand{\Pmodel}{P_{\rm{model}}}
\newcommand{\ptildemodel}{\tilde{p}_{\rm{model}}}
% Stochastic autoencoder distributions
\newcommand{\pencode}{p_{\rm{encoder}}}
\newcommand{\pdecode}{p_{\rm{decoder}}}
\newcommand{\precons}{p_{\rm{reconstruct}}}

\newcommand{\laplace}{\mathrm{Laplace}} % Laplace distribution

\newcommand{\E}{\mathbb{E}}
\newcommand{\Ls}{\mathcal{L}}
\newcommand{\R}{\mathbb{R}}
\newcommand{\emp}{\tilde{p}}
\newcommand{\lr}{\alpha}
\newcommand{\reg}{\lambda}
\newcommand{\rect}{\mathrm{rectifier}}
\newcommand{\softmax}{\mathrm{softmax}}
\newcommand{\sigmoid}{\sigma}
\newcommand{\softplus}{\zeta}
\newcommand{\KL}{D_{\mathrm{KL}}}
\newcommand{\Var}{\mathrm{Var}}
\newcommand{\standarderror}{\mathrm{SE}}
\newcommand{\Cov}{\mathrm{Cov}}
% Wolfram Mathworld says $L^2$ is for function spaces and $\ell^2$ is for vectors
% But then they seem to use $L^2$ for vectors throughout the site, and so does
% wikipedia.
\newcommand{\normlzero}{L^0}
\newcommand{\normlone}{L^1}
\newcommand{\normltwo}{L^2}
\newcommand{\normlp}{L^p}
\newcommand{\normmax}{L^\infty}

\newcommand{\parents}{Pa} % See usage in notation.tex. Chosen to match Daphne's book.

\DeclareMathOperator*{\argmax}{arg\,max}
\DeclareMathOperator*{\argmin}{arg\,min}

\DeclareMathOperator{\sign}{sign}
\DeclareMathOperator{\Tr}{Tr}
\let\ab\allowbreak


\newtheorem{theorem}{Theorem}
\newtheorem{lemma}{Lemma}
\newtheorem{assumption}{Assumption}
\newtheorem{example}{Example}
\newtheorem{proposition}{Proposition}
\newtheorem{corollary}{Corollary}
\newtheorem{remark}{Remark}

\title{Stochastic Compositional Optimization via \\ Hybrid Momentum Frank--Wolfe}

\author{%
  El Mahdi Chayti\\
  Machine Learning and Optimization Laboratory (MLO)\\
  EPFL, Switzerland\\
  \texttt{el-mahdi.chayti[AT]epfl.ch}
}

\begin{document}

\maketitle

\begin{abstract}
Stochastic compositional optimization minimizes objectives of the form $\min_{\vx \in \cX} F(\vf(\vx), \vx)$, where $\vf$ is accessible only through noisy stochastic queries. Existing methods for this problem assume that the outer function $F$ is continuously differentiable, which excludes many practically important applications such as robust max-of-losses, Conditional Value-at-Risk, and norm regularizers. We propose the Hybrid Momentum Stochastic Frank--Wolfe algorithm, which drops the smoothness assumption on $F$. By combining a momentum-based Jacobian tracker with a Taylor-corrected function tracker, the algorithm feeds an entire stochastic linearization---rather than a single gradient---into a generalized linear minimization oracle. We establish an $\mathcal{O}(K^{-1/4})$ convergence rate in the generalized Frank--Wolfe gap for non-convex objectives with $L_F$-Lipschitz outer functions, matching the optimal complexity for projection-free single-sample stochastic methods under expected smoothness. The analysis extends to heavy-tailed noise oracles with bounded $r$-th moments for $r \in (1, 2]$ and recovers the deterministic rates of \citet{vladarean2023first} as the noise vanishes.
\end{abstract}

\section{Introduction}

Modern machine learning increasingly demands optimization paradigms that go beyond classical Empirical Risk Minimization. Risk-averse learning, distributionally robust optimization, multi-task learning, and fairness-constrained training all give rise to objectives of the form
\begin{equation}
    \min_{\vx \in \cX} \varphi(\vx) := F(\vf(\vx), \vx) \quad \text{where} \quad \vf(\vx) = \E_{\xi}[\tilde{\vf}(\vx; \xi)], \label{eq:problem}
\end{equation}
where $\cX \subset \R^d$ is convex and compact, $\vf : \cX \to \R^n$ is a vector-valued inner mapping accessible only via stochastic queries, and $F : \R^n \times \cX \to \R$ is a deterministic outer function. This is the \emph{fully composite} formulation introduced in  \citet{doikov2022high} and revisited in its projection-free variant by \citep{vladarean2023first} (in the deterministic case); this formulation generalizes both classical compositional optimization $F(\vf(\vx))$ and additive composite optimization.

Solving \eqref{eq:problem} stochastically poses a fundamental challenge: by Jensen's inequality, naive plug-in estimators are biased, $\E[F(\tilde\vf(\vx; \xi), \vx)] \ne F(\E[\tilde\vf(\vx; \xi)], \vx)$. A line of work culminating in \citet{chen2020solving, balasubramanian2022stochastic, xiao2022projection} addresses this by tracking the inner function with momentum-based estimators and feeding the result through the chain rule $\nabla \varphi(\vx) = \nabla \vf(\vx)^\top \nabla F(\vf(\vx))$. However, the chain rule \emph{requires $F$ to be differentiable}---an assumption that excludes the very examples that motivate fully composite optimization in the first place: max-type losses, CVaR, $\ell_1$ regularizers, and constraint-penalty formulations.

\citet{vladarean2023first} resolved this in the deterministic setting through a \emph{generalized linear minimization oracle} (GLMO), which minimizes $F$ evaluated on a linearization of $\vf$, treating $F$ as a black box rather than differentiating through it. They explicitly leave the stochastic case as an open problem.

\paragraph{Contributions.} This paper addresses that open problem.
\begin{enumerate}
\item We propose the \emph{Hybrid Momentum Stochastic Frank--Wolfe} algorithm (Algorithm~\ref{alg:hybrid_sfw}), the first stochastic projection-free method for fully composite problems with non-smooth outer $F$. The algorithm tracks both the Jacobian and the function value of $\vf$, and feeds the resulting affine model into a GLMO that handles non-smooth $F$ directly.
\item We present two tracker variants offering a memory/assumption trade-off: \emph{Variant I} (Polyak update; memoryless but requires $\vf$ Lipschitz); \emph{Variant II} (Taylor-corrected; requires storing one previous iterate but eliminates the Lipschitz assumption).
\item Under heavy-tailed noise with bounded $r$-th moment ($r \in (1, 2]$), we establish $\mathcal{O}(K^{-(r-1)/(3r-2)})$ convergence in the generalized Frank--Wolfe gap (Theorem~\ref{thm:nonconvex}). At $r = 2$, this gives $\mathcal{O}(K^{-1/4})$, matching the minimax-optimal rate for single-sample projection-free stochastic methods under expected smoothness \citep{arjevani2023lower}.
\item For convex objectives, we obtain $\mathcal{O}(K^{-(r-1)/(2r-1)})$ primal rate (Theorem~\ref{thm:convex}), recovering $\mathcal{O}(K^{-1/3})$ at $r = 2$.
\item The algorithm interpolates with the deterministic regime: when the oracle is exact, setting the momentum parameters to one recovers the deterministic Basic Method of \citet{vladarean2023first} together with its $\mathcal{O}(1/K)$ convex and $\tilde{\mathcal{O}}(1/\sqrt{K})$ non-convex rates (Corollary~\ref{cor:deterministic}).
\item All constants are derived explicitly in terms of problem parameters, and we provide numerical experiments on three real-world problems exhibiting heavy-tailed noise: robust minimax sparse regression on \texttt{a9a} \citep{chang2011libsvm}, CVaR portfolio optimization on S\&P~500 returns \citep{wrds_compustat}, and robust matrix completion on \texttt{MovieLens-100K} \citep{harper2015movielens}.
\end{enumerate}

\paragraph{Related work.}
The deterministic foundation for our approach was established by \citet{vladarean2023first}, who introduced the GLMO and analyzed the deterministic Basic and Accelerated methods for fully composite optimization. In a related vein, \citet{doikov2022high} investigated higher-order methods for this same problem class.

In the stochastic regime, early work by \citet{wang2017stochastic} introduced SCGD using a two-timescale stepsize schedule. This was later streamlined by \citet{ghadimi2020single} via the single-timescale NASA algorithm. Subsequent advancements include SCSC \citep{chen2020solving}, which achieved optimal SGD-like rates, and Linearized NASA \citep{balasubramanian2022stochastic}, which established level-independent rates. However, all of these methods rely on exact projections and cannot operate in the projection-free setting.

For projection-free stochastic compositional optimization, the current state-of-the-art is LiNASA+ICG \citep{xiao2022projection}. While it achieves an $\mathcal{O}(K^{-1/4})$ convergence rate, it suffers from a high per-iteration cost, requiring $\mathcal{O}(K^{3/2})$ LMO calls (yielding an overall $\mathcal{O}(\epsilon^{-4})$ SFO and $\mathcal{O}(\epsilon^{-6})$ LMO complexity). Furthermore, it fundamentally requires the outer function $F$ to be smooth. Similarly, \citet{akhtar2021projection} tackle the two-level case but impose strictly stronger assumptions. Crucially, both of these methods require differentiating through $F$, rendering them inapplicable to the non-smooth outer functions considered in our work.

Finally, our variance reduction techniques draw inspiration from the standard (non-compositional) stochastic Frank--Wolfe literature. \citet{mokhtari2020stochastic} demonstrated an $\mathcal{O}(K^{-1/3})$ rate for convex objectives using momentum tracking, while \citet{yurtsever2019conditional} achieved improved non-convex rates utilizing the STORM variance reduction technique originally developed by \citet{cutkosky2019momentum}. To correct the bias inherent in momentum tracking, \citet{hassani2020stochastic} and \citet{zhang2019samplestochasticfrankwolfe} proposed using Hessian-vector products coupled with uniform random variables. This bias correction was recently refined by \citet{chayti2026ransomsecondordermomentumrandomized} using Beta-distributed random variables. We adapt both of these fundamental techniques to our fully composite framework, detailing a STORM-based acceleration in Appendix~\ref{sec:storm_extension} and a Hessian-based acceleration in Appendix~\ref{sec:hessian_extension}.

\section{Problem Setup and Assumptions}

\paragraph{Notation.} We use $\|\cdot\|$ for the Euclidean norm (Frobenius norm for matrices), $D_{\cX} := \max_{\vx, \vy \in \cX} \|\vx - \vy\|$ for the diameter of $\cX$, and $\mathbf{1} \in \R^n$ for the all-ones vector. 

\textbf{Generalized Frank--Wolfe gap.} For $\vy \in \cX$, the \emph{generalized Frank--Wolfe gap} is
\begin{equation} \label{eq:gen_fw_gap}
\hat\Delta(\vy) := \varphi(\vy) - \min_{\vx \in \cX} F\big(\vf(\vy) + \nabla \vf(\vy)(\vx - \vy),\, \vx\big) \ge 0,
\end{equation}
which generalizes the classical FW gap i.e.: when $F(\vu, \vx) \equiv u^{(1)}$, $\hat\Delta(\vy)$ reduces to $\max_{\vx \in \cX} \langle \nabla f_1(\vy), \vy - \vx \rangle$.

Note that by construction, $\hat\Delta(\vy) \ge 0$, and $\hat\Delta(\vy) = 0$ characterizes stationarity in the non-smooth setting \citep{vladarean2023first}.

\subsection{Assumptions}

\begin{assumption}[Domain compactness] \label{assump:domain}
$\cX \subset \R^d$ is convex and compact with diameter $D_{\cX} < \infty$.
\end{assumption}

\begin{assumption}[Outer function regularity] \label{assump:outer}
$F : \R^n \times \cX \to \R$ is (i) $L_F$-Lipschitz in $\vu$, (ii) monotone non-decreasing in $\vu$ (i.e., $\vu_1 \le \vu_2$ componentwise implies $F(\vu_1, \vx) \le F(\vu_2, \vx)$), and (iii) subhomogeneous in $\vu$: $F(\gamma \vu, \vx) \le \gamma F(\vu, \vx)$ for all $\gamma \ge 1$. $F$ is \emph{not} assumed to be differentiable.
\end{assumption}

\begin{assumption}[Inner function smoothness] \label{assump:inner_smooth}
$\vf : \cX \to \R^n$ has $L$-Lipschitz Jacobian: $\|\nabla \vf(\vx) - \nabla \vf(\vy)\|_{\mathrm{op}} \le L \|\vx - \vy\|$ for all $\vx, \vy \in \cX$.
\end{assumption}

\begin{assumption}[Heavy-tailed stochastic oracle] \label{assump:noise}
There exists $r \in (1, 2]$ and constants $\sigma_f, \sigma_g \ge 0$ such that for every $\vx \in \cX$, the stochastic oracle returns unbiased estimates $\tilde\vf(\vx; \xi), \nabla\tilde\vf(\vx; \xi)$ satisfying
\[
\E_\xi\|\tilde\vf(\vx; \xi) - \vf(\vx)\|^r \le \sigma_f^r,
\qquad
\E_\xi\|\nabla \tilde\vf(\vx; \xi) - \nabla \vf(\vx)\|^r \le \sigma_g^r.
\]
The classical bounded-variance setting corresponds to $r = 2$.
\end{assumption}

The following two assumptions are used in specific results.

\begin{assumption}[Lipschitz inner function] \label{assump:inner_lipschitz}
$\vf$ is $G$-Lipschitz: $\|\vf(\vx) - \vf(\vy)\| \le G \|\vx - \vy\|$ for all $\vx, \vy \in \cX$. Used only for Variant I.
\end{assumption}

\begin{assumption}[Convexity for the convex setting] \label{assump:convex}
Each $f_i$ is convex on $\cX$ and $F(\vu, \vx)$ is jointly convex in $(\vu, \vx)$. Used only for Theorem~\ref{thm:convex}; combined with monotonicity of $F$, this implies $\varphi$ is convex on $\cX$.
\end{assumption}

\paragraph{Discussion of the assumptions.}
Assumptions~\ref{assump:domain}, \ref{assump:outer}(ii)--(iii), \ref{assump:inner_smooth}, and \ref{assump:convex} are standard in the deterministic fully composite optimization literature \citep{vladarean2023first, doikov2022high}. Assumption~\ref{assump:noise} is standard in stochastic optimization with heavy-tailed gradients and reduces to bounded variance at $r = 2$. Assumption~\ref{assump:inner_lipschitz} is standard for stochastic gradient methods.

The single new condition in our setting is Assumption~\ref{assump:outer}(i), the Lipschitzness of $F$ in $\vu$. This is necessary because $F$ is non-smooth: in the absence of differentiability, Lipschitzness is the natural quantitative substitute that controls the deviation between the GLMO evaluated on the exact and tracked surrogates (Lemma~\ref{lem:glmo_main}). The Lipschitz constant $L_F$ admits closed-form expressions for all the motivating examples below: $L_F = 1$ for max-of-losses, $\ell_p$-norm regularization, and composite regularization; $L_F \le 1/(\alpha\sqrt{n})$ for CVaR at level $\alpha$; and $L_F \le \sqrt{(1+\rho)/n}$ for $\chi^2$-DRO with radius $\rho$.

\subsection{Composite Curvature}

The composite curvature constant generalizes the standard FW curvature to the non-smooth outer setting:
\begin{equation} \label{eq:curvature_def}
\mathcal{S} := \sup_{\substack{\vx, \vy \in \cX \\ \gamma \in (0, 1]}} \frac{2}{\gamma^2} \Big[ F(\vf(\vy_\gamma), \vy_\gamma) - F(\vf(\vx) + \gamma \nabla \vf(\vx)(\vy - \vx), \vy_\gamma) \Big],
\end{equation}
where $\vy_\gamma := \vx + \gamma(\vy - \vx)$. Under Assumptions~\ref{assump:domain}--\ref{assump:inner_smooth}, $\mathcal{S} \le L_F L D_{\cX}^2 \sqrt{n}$ (Proposition~\ref{prop:curvature} in the appendix).

\section{Algorithm}

The algorithm maintains two stochastic trackers: a Polyak Jacobian tracker $\mV_k$ and a function-value tracker $\vz_k$ in one of two variants. At each step, the trackers define an affine surrogate $\tilde l_k(\vx) := \vz_k + \mV_k(\vx - \vy_k)$, fed into the GLMO.

\begin{algorithm}[h]
\caption{Hybrid Momentum Stochastic Frank--Wolfe}
\label{alg:hybrid_sfw}
\begin{algorithmic}[1]
\STATE \textbf{Input:} $\vy_0 \in \cX$, step-size sequences $\{\gamma_k, \beta_k, \rho_k\} \subset (0, 1]$, total iterations $K$.
\STATE Draw $\xi_0$; set $\vz_0 \!=\! \tilde\vf(\vy_0; \xi_0)$, $\mV_0 \!=\! \nabla \tilde\vf(\vy_0; \xi_0)$.
\FOR{$k = 1, \ldots, K$}
\STATE Draw $\xi_k$ at $\vy_k$.
\STATE \textbf{Jacobian tracker:} $\mV_k = (1-\beta_k)\mV_{k-1} + \beta_k \nabla \tilde\vf(\vy_k; \xi_k)$.
\STATE \textbf{Function tracker} (one variant):
\begin{align*}
\text{Variant I (Polyak):}\quad & \vz_k = (1-\rho_k)\vz_{k-1} + \rho_k \tilde\vf(\vy_k; \xi_k), \\
\text{Variant II (Taylor):}\quad & \vz_k = (1-\rho_k)[\vz_{k-1} + \mV_k(\vy_k - \vy_{k-1})] + \rho_k \tilde\vf(\vy_k; \xi_k).
\end{align*}
\STATE \textbf{GLMO:} $\vx_{k+1} \in \argmin_{\vx \in \cX} F(\vz_k + \mV_k(\vx - \vy_k),\, \vx)$.
\STATE \textbf{Update:} $\vy_{k+1} = (1-\gamma_k)\vy_k + \gamma_k \vx_{k+1}$.
\ENDFOR
\STATE \textbf{Return:} $\vy_K$.
\end{algorithmic}
\end{algorithm}

\paragraph{Variant trade-off.}
Variant I is memoryless and requires only Lipschitz $\vf$ (Assumption~\ref{assump:inner_lipschitz}); its tracking error scales linearly with the step distance. Variant II uses the first-order Taylor extrapolation $\mV_k(\vy_k - \vy_{k-1})$ to predict $\vf(\vy_k)$ from $\vf(\vy_{k-1})$, reducing the tracking error to a quadratic Taylor remainder; it eliminates Assumption~\ref{assump:inner_lipschitz} but requires storing the previous iterate.

\paragraph{Generalized LMO.}
The GLMO at line~7 minimizes $F$ evaluated on the affine surrogate $\tilde l_k(\vx)$. As shown by \citet{vladarean2023first}, this oracle is tractable for all the motivating examples below: it reduces to a small linear program (or, in the case of $\ell_2$-norm or $\chi^2$-DRO, a second-order cone program). Crucially, the GLMO handles the non-smoothness of $F$ directly, without requiring a (sub)gradient.
We also note that the GLMO reduces to the classical LMO used in Frank--Wolfe methods \citep{frank1956algorithm,jaggi2013revisiting} when $F(\vu, \vx) \equiv u^{(1)}$.

\section{Motivating Examples}

We illustrate the GLMO on five paradigmatic non-smooth fully composite problems. In each case, the stochastic inner function is $\vf(\vx) = \E_\xi[\tilde\vf(\vx; \xi)]$.

\begin{example}[Robust Minimax Optimization]
Consider $n$ distinct competing objectives or losses, $f_i(\vx) = \E_{\xi_i}[\tilde{f}_i(\vx; \xi_i)]$. We aim to minimize the worst-case expected loss over a structured domain (e.g., an $\ell_1$-ball for sparsity) \citep{shalev2016minimizing}:
\begin{equation}
    \min_{\vx \in \cX} \max_{1 \le i \le n} f_i(\vx).
\end{equation}
Here, the outer function is $F(\vu) = \max_i u^{(i)}$, which is globally $1$-Lipschitz but non-differentiable. \textbf{Tractability:} The GLMO requires minimizing the maximum of $n$ affine functions over $\cX$, which trivially reformulates as a highly efficient Linear Program (LP) by introducing a single auxiliary scalar variable \citep{vladarean2023first}.
\end{example}


\begin{example}[Norm-Regularized Multi-Task Learning]
In multi-task learning, one often minimizes the norm of a vector of expected tasks to encourage joint sparsity or uniformly bounded errors across tasks \citep{sener2018multi}:
\begin{equation}
    \min_{\vx \in \cX} \|\vf(\vx)\|_p.
\end{equation}
Here, $F(\vu) = \|\vu\|_p$. For $p \ge 1$, $F$ is globally Lipschitz but non-smooth at the origin. \textbf{Tractability:} For $p=1$ or $\infty$, the GLMO reduces to an LP. For $p=2$, it translates to an SOCP, efficiently solvable via interior point methods.
\end{example}

\begin{example}[Conditional Value at Risk (CVaR)]
Risk-averse learning often minimizes the CVaR at level $\alpha \in (0,1)$, which focuses on the $\alpha$-fraction of the worst-case expected losses across $n$ tasks \citep{rockafellar2000optimization}:
\begin{equation}
    \min_{\vx \in \cX} \text{CVaR}_\alpha(\vf(\vx)) = \min_{\vx \in \cX} \inf_{\tau \in \R} \left\{ \tau + \frac{1}{\alpha n} \sum_{i=1}^n \max(0, f_i(\vx) - \tau) \right\}.
\end{equation}
Here, $F(\vu)$ encapsulates the CVaR operator. \textbf{Tractability:} Evaluating the GLMO simply requires minimizing a piecewise affine convex function over $\cX$. By adding auxiliary variables for the rectified linear units ($\max(0, \cdot)$), the GLMO cleanly formulates as an LP \citep{vladarean2023first}.
\end{example}

\begin{example}[Distributionally Robust Optimization (DRO)]
In DRO, one minimizes the expected loss with respect to the worst-case distribution $\mathbf{q}$ within a $\chi^2$-divergence ball of radius $\rho$ around the uniform distribution \citep{namkoong2016stochastic}:
\begin{equation}
    \min_{\vx \in \cX} \max_{\mathbf{q} \in \Delta_n, D_{\chi^2}(\mathbf{q}, \frac{1}{n}\mathbf{1}) \le \rho} \sum_{i=1}^n q_i f_i(\vx).
\end{equation}
Here, $F(\vu)$ is the supremum over the $\chi^2$ ball. \textbf{Tractability:} Through convex duality, this minimax objective can be rewritten as a minimization over an affine surrogate, reducing the GLMO to a low-dimensional convex problem bounded by a quadratic constraint (an SOCP).
\end{example}

\begin{example}[Composite Regularized Optimization]
Consider a standard expected risk minimization problem augmented with a deterministic, potentially non-smooth, convex regularizer $\Xi(\vx)$ \citep{yurtsever2018conditional}:
\begin{equation}
    \min_{\vx \in \cX} \E_{\xi}[\tilde{f}(\vx; \xi)] + \Xi(\vx).
\end{equation}
Here, $f(\vx) = \E[\tilde{f}(\vx)]$, and the outer function is explicitly dependent on $\vx$: $F(u, \vx) = u + \Xi(\vx)$. \textbf{Tractability:} The GLMO computes $\argmin_{\vx \in \cX} \left[ \vz_k + \langle \mV_k, \vx - \vy_k \rangle + \Xi(\vx) \right]$. Discarding constant terms, this strictly reduces to $\argmin_{\vx \in \cX} \langle \mV_k, \vx \rangle + \Xi(\vx)$, which exactly recovers the highly tractable classical composite Frank--Wolfe oracle \citep{yurtsever2018conditional}.
\end{example}

Further examples covering fairness-constrained learning \citep{agarwal2018reductions, donini2018empirical} and exact penalty methods \citep{cotter2019two} are deferred to Appendix~\ref{sec:examples_appendix}.

\section{Convergence Analysis}

We present the main convergence guarantees for non-convex objectives (Theorem~\ref{thm:nonconvex}) and convex objectives (Theorem~\ref{thm:convex}), and discuss the deterministic-noiseless interpolation (Corollary~\ref{cor:deterministic}). All proofs are deferred to Appendices~\ref{sec:proof_nonconvex} and \ref{sec:proof_convex}.

\subsection{Non-Convex Convergence}

\begin{theorem}[Non-convex convergence] \label{thm:nonconvex}
Suppose Assumptions~\ref{assump:domain}--\ref{assump:noise} hold. Choose the constant step sizes
\begin{equation}\label{eq:nonconvex_schedule}
\gamma_k \equiv \gamma = K^{-(2r-1)/(3r-2)}, \qquad \beta_k = \rho_k \equiv \beta = K^{-r/(3r-2)}.
\end{equation}
Then for each variant, Algorithm~\ref{alg:hybrid_sfw} satisfies
\[
\min_{1 \le k \le K} \E[\hat\Delta_k] = \mathcal{O}\!\left(K^{-(r-1)/(3r-2)}\right).
\]
Variant I requires the additional Assumption~\ref{assump:inner_lipschitz}; Variant II does not. Explicit constants $M_{\mathrm{I}}$ and $M_{\mathrm{II}}$ in the $\mathcal{O}(\cdot)$ notation are given in equations \eqref{eq:MI_explicit}--\eqref{eq:MII_explicit} of Appendix~\ref{sec:proof_nonconvex}.
\end{theorem}

\paragraph{Proof sketch.} The proof combines a one-step progress bound for the GLMO update (Lemma~\ref{lem:glmo_main}) with bounds on the $r$-th moment of the tracking errors (Lemmas~\ref{lem:grad_tracker_main}--\ref{lem:func_taylor_main} for the Jacobian, Polyak, and Taylor trackers respectively). Telescoping, applying Jensen's inequality, and balancing terms with the schedule \eqref{eq:nonconvex_schedule} yields the rate. See Appendix~\ref{sec:proof_nonconvex}.

\paragraph{Optimality.}
At $r = 2$, Theorem~\ref{thm:nonconvex} gives the rate $\mathcal{O}(K^{-1/4})$, matching the lower bound of \citet{arjevani2023lower} for single-sample stochastic methods under expected smoothness. Faster rates such as $\mathcal{O}(K^{-1/3})$ via STORM-style momentum \citep{cutkosky2019momentum} or via second-order corrections \citep{hassani2020stochastic} require strictly stronger assumptions. In particular, Proposition~\ref{prop:hessian_implies_avg} (Appendix~\ref{sec:hessian_extension}) shows that bounded Hessian-noise moments imply $r$-average smoothness, so any acceleration to the $K^{-(r-1)/(2r-1)}$ rate is fundamentally tied to a stronger smoothness condition.

\subsection{Convex Convergence}

\begin{theorem}[Convex convergence] \label{thm:convex}
Suppose Assumptions~\ref{assump:domain}--\ref{assump:noise} and \ref{assump:convex} hold. Let $\vx^* \in \argmin_{\vx \in \cX} \varphi(\vx)$ and $\Phi_0 := \varphi(\vy_0) - \varphi(\vx^*)$. Choose the decreasing step sizes
\begin{equation}\label{eq:convex_schedule}
\gamma_k = \frac{2}{k+2}, \qquad \beta_k = \rho_k = \min\!\left\{1,\ \frac{c_0}{(k+k_0)^{r/(2r-1)}}\right\},
\end{equation}
with $c_0$ and $k_0$ specified in Appendix~\ref{sec:proof_convex}. Then for each variant of Algorithm~\ref{alg:hybrid_sfw} (Variant I additionally assuming Assumption~\ref{assump:inner_lipschitz}),
\[
\E[\varphi(\vy_K) - \varphi(\vx^*)] \;\le\; \E[\hat\Delta_K] \;=\; \mathcal{O}\!\left((K+2)^{-(r-1)/(2r-1)}\right).
\]
The explicit constant $A_{\mathrm{cvx}}$ is given in equation~\eqref{eq:Acvx_def} of Appendix~\ref{sec:proof_convex}.
\end{theorem}

At $r = 2$, this yields $\mathcal{O}(K^{-1/3})$, matching the optimal rate for convex stochastic FW with momentum \citep{mokhtari2020stochastic}, now extended to the non-smooth fully composite setting.

\subsection{Deterministic Interpolation and Nonconvex Acceleration}

A natural sanity check is that the algorithm reduces to the deterministic Frank--Wolfe of \citet{vladarean2023first} when the oracle is exact. The following corollary makes this precise.

\begin{corollary}[Deterministic interpolation] \label{cor:deterministic}
Suppose the stochastic oracle is exact ($\sigma_f = \sigma_g = 0$ in Assumption~\ref{assump:noise}). Setting $\beta_k = \rho_k = 1$ for all $k$ in Algorithm~\ref{alg:hybrid_sfw} makes the trackers exact: $\mV_k = \nabla \vf(\vy_k)$ and $\vz_k = \vf(\vy_k)$ for all $k \ge 1$. The algorithm then reduces to the Basic Method of \citet{vladarean2023first}. With the deterministic step-size schedules, the iterates satisfy:
\begin{itemize}
\item \textbf{Convex case} (Assumption~\ref{assump:convex}, $\gamma_k = 2/(k+2)$):
\[
\varphi(\vy_K) - \varphi(\vx^*) \le \frac{2 \mathcal{S}}{K+1}.
\]
\item \textbf{Non-convex case} ($\gamma_k = 1/\sqrt{k+1}$):
\[
\min_{0 \le k \le K} \hat\Delta_k \le \frac{\Phi_0 + \tfrac{1}{2}\mathcal{S}(1 + \ln(K+1))}{\sqrt{K+1}}.
\]
\end{itemize}
These match Theorems 3.1 and 3.2 of \citet{vladarean2023first}.
\end{corollary}

The corollary is immediate from the form of the algorithm and Lemma~\ref{lem:glmo_main}: with exact trackers, the tracking-error contributions vanish and Lemma~\ref{lem:glmo_main} reduces to the deterministic progress bound. See Appendix~\ref{sec:proof_deterministic} for details.

\paragraph{Schedule choice across regimes.}
The schedules in Theorems~\ref{thm:nonconvex} and~\ref{thm:convex} are noise-aware: they shrink the momentum parameters $\beta_k, \rho_k$ as $K$ grows in order to absorb the variance of the stochastic estimates. When $\sigma_f = \sigma_g = 0$, no shrinking is needed: setting $\beta_k = \rho_k = 1$ recovers the deterministic regime, where Vladarean's $\gamma_k$-only schedule is optimal. The algorithm thus interpolates smoothly: the same code, with three momentum parameters, recovers either the stochastic or the deterministic algorithm depending on how those parameters are set.

\paragraph{Acceleration under stronger smoothness.}
The rates in Theorems~\ref{thm:nonconvex} and~\ref{thm:convex} are tight under expected smoothness alone (Assumption~\ref{assump:inner_smooth}). When the stochastic Jacobian satisfies a stronger \emph{$r$-average smoothness} property (Assumption~\ref{assump:avg_smooth} in Appendix~\ref{sec:storm_extension}), our framework integrates seamlessly with the STORM variance-reduction technique \citep{cutkosky2019momentum, hassani2020stochastic}. Specifically, replacing the Polyak Jacobian update by the STORM finite-difference correction (\eqref{eq:storm_grad} in the appendix), with the same Taylor-corrected function tracker, accelerates the rate to $\mathcal{O}(K^{-(r-1)/(2r-1)})$, recovering $\mathcal{O}(K^{-1/3})$ at $r = 2$ (Theorem~\ref{thm:storm_rate}). An alternative second-order acceleration via randomized Hessian--vector products is given in Appendix~\ref{sec:hessian_extension}; we additionally prove there (Proposition~\ref{prop:hessian_implies_avg}) that bounded-Hessian-noise moments imply $r$-average smoothness, showing that all such accelerations are tied to the same underlying structural assumption.

\subsection{Summary of Results}

Table~\ref{tab:comparison} situates our results within the existing literature on stochastic compositional optimization. Our work is the first to handle a non-smooth outer function $F$ in the stochastic projection-free regime. Among prior work, only the deterministic Basic Method of \citet{vladarean2023first} permits a non-smooth $F$, and only \citet{xiao2022projection} operates projection-free in the stochastic case (though it strictly requires $F$ to be smooth).

\begin{table}[h]
\centering
\small
\caption{Comparison of stochastic compositional optimization methods. Rates are stated for finite-variance noise ($r = 2$); extensions to heavy-tailed noise with $r \in (1, 2]$ are given in the cited references where applicable. ``PF'' denotes projection-free; ``Smooth'' refers to smoothness of the outer function $F$; ``Inner'' denotes the regularity assumption on the stochastic Jacobian (Exp.\ = expected smoothness, Avg.\ = $r$-average smoothness).}
\label{tab:comparison}
\begin{tabular}{@{}lccccll@{}}
\toprule
Method & Stoch. & PF & Smooth $F$ & Inner & Setting & Rate \\
\midrule
SCGD \citep{wang2017stochastic}        & \checkmark & --        & required & Exp. & nonconvex & $\mathcal{O}(K^{-1/4})$ \\
NASA \citep{ghadimi2020single}         & \checkmark & --        & required & Exp. & nonconvex & $\mathcal{O}(K^{-1/4})$ \\
SCSC \citep{chen2020solving}           & \checkmark & --        & required & Exp. & nonconvex & $\mathcal{O}(K^{-1/4})$ \\
L-NASA \citep{balasubramanian2022stochastic} & \checkmark & --  & required & Exp. & nonconvex & $\mathcal{O}(K^{-1/4})$ \\
LiNASA+ICG \citep{xiao2022projection}  & \checkmark & \checkmark & required & Exp. & nonconvex & $\mathcal{O}(K^{-1/4})$ \\
Vladarean et al.\ \citep{vladarean2023first} & --      & \checkmark & \emph{not required} & --   & nonconvex & $\tilde{\mathcal{O}}(K^{-1/2})$ \\
Vladarean et al.\ \citep{vladarean2023first} & --      & \checkmark & \emph{not required} & --   & convex    & $\mathcal{O}(K^{-1})$ \\
\midrule
\textbf{This work} (Thm.~\ref{thm:nonconvex})  & \checkmark & \checkmark & \emph{not required} & Exp. & nonconvex & $\mathcal{O}(K^{-1/4})$ \\
\textbf{This work} (Thm.~\ref{thm:convex})     & \checkmark & \checkmark & \emph{not required} & Exp. & convex    & $\mathcal{O}(K^{-1/3})$ \\
\textbf{This work} (Thm.~\ref{thm:storm_rate})& \checkmark & \checkmark & \emph{not required} & Avg. & nonconvex & $\mathcal{O}(K^{-1/3})$ \\
\bottomrule
\end{tabular}
\end{table}

Four observations stand out. First, our method is the only stochastic projection-free algorithm that does not require $F$ to be smooth, addressing the open problem stated in \citet{vladarean2023first}. Second, the $\mathcal{O}(K^{-1/4})$ rate of Theorem~\ref{thm:nonconvex} matches the established state-of-the-art for projection-based stochastic compositional methods such as SCSC and L-NASA (which report $\mathcal{O}(K^{-1/2})$ bounds evaluated with respect to the \emph{squared} gap). Furthermore, this $\mathcal{O}(K^{-1/4})$ rate is minimax-optimal for projection-free single-sample stochastic methods under expected smoothness, as established by the lower bounds of \citet{arjevani2023lower}. Third, regarding oracle complexity, our framework requires only a single Generalized Linear Minimization Oracle (GLMO) call per iteration. In contrast, the only other stochastic projection-free method, LiNASA+ICG \citep{xiao2022projection}, relies on an inner Inverse Conditional Gradient (ICG) procedure that demands $\mathcal{O}(K^{3/2})$ standard LMO calls to achieve its $\mathcal{O}(K^{-1/4})$ rate, making our approach significantly more computationally efficient. Fourth, when the stronger $r$-average smoothness assumption is available, Theorem~\ref{thm:storm_rate} accelerates this rate to $\mathcal{O}(K^{-1/3})$ via STORM-style variance reduction. This recovers the accelerated rate of \citet{yurtsever2019conditional} for standard stochastic Frank--Wolfe, but successfully extends it to the much broader class of problems featuring non-smooth compositional objectives.

\section{Numerical Experiments} \label{sec:experiments}

We evaluate the algorithm on three problems featuring non-smooth $F$ and non-trivial constraint sets: robust minimax sparse regression, CVaR portfolio optimization, and robust low-rank matrix completion.

\textbf{Algorithms.}
We benchmark four algorithms, each using the same GLMO. (i) \textbf{Vanilla SCFW}: stochastic compositional FW with no variance reduction. (ii) \textbf{Clipped SCFW}: Vanilla SCFW with the stochastic Jacobian clipped at norm $C$ (tuned per task). (iii) \textbf{Variant I} (this work, Polyak tracker). (iv) \textbf{Variant II} (this work, Taylor tracker). For Variants I and II, we use the theoretically-prescribed schedule at $r = 2$: $\gamma = K^{-3/4}$, $\beta = \rho = K^{-1/2}$. We evaluate progress via the generalized FW gap $\hat\Delta_k$, plotted on log-log scale.

\paragraph{Tasks.}
We consider the three following tasks:
\begin{enumerate}[label=\textbf{Task \arabic*:}, leftmargin=*,noitemsep,topsep=0pt]
    \item \textbf{Robust minimax sparse regression.} Given $m$ data subsets $\{(\va_{i,j}, b_{i,j})\}_{j=1}^{N_i}$, $i = 1, \ldots, m$, we solve
    \begin{equation}
        \min_{\|\vx\|_1 \le \tau} \max_{1 \le i \le m} \tfrac{1}{2 N_i} \sum_{j=1}^{N_i} (\va_{i,j}^\top \vx - b_{i,j})^2,
    \end{equation}
    where the outer function $F(\vu) = \max_i u^{(i)}$ is non-smooth. The constraint $\cX = \{\vx : \|\vx\|_1 \le \tau\}$ admits an $\mathcal{O}(d)$ LMO, and the GLMO is a small linear program in $\vx$ and an auxiliary scalar.

    \item \textbf{Risk-averse portfolio (CVaR).} Given asset return vectors $\vr_t \in \R^d$, we solve
    \begin{equation}
        \min_{[\vx, y_0] \in \Delta_d \times [-1, 1]} \Big\{ y_0 + \frac{1}{1 - \alpha} \E_{\vr}\big[\max(0, -\vr^\top \vx - y_0)\big] \Big\},
    \end{equation}
    where $\Delta_d$ is the probability simplex and $\alpha = 0.95$. The outer function involves the non-smooth hinge operator. The GLMO is an LP.

    \item \textbf{Robust matrix completion.} Given partial observations $\Omega \subset [d] \times [m]$ of a target matrix $\mM$, we solve
    \begin{equation}
        \min_{\|\mX\|_* \le \tau} \tfrac{1}{|\Omega|} \sum_{(i,j) \in \Omega} |X_{ij} - M_{ij}|,
    \end{equation}
    where the $\ell_1$ loss $F(\vu) = \|\vu\|_1$ is non-smooth and the nuclear-norm ball admits an LMO via top-singular-vector computation.
\end{enumerate}

\textbf{Synthetic-data results.}
We first verify the convergence behavior on synthetic instances where ground-truth quantities are known and the noise is controlled. Task~1 uses $m = 10$ data subsets with $d = 100$, $\tau = 5$, and Gaussian inputs with task-dependent noise. Task~2 uses $d = 50$ assets with synthetically generated heavy-tailed daily returns. Task~3 uses a $30 \times 20$ matrix with rank $5$ and $30\%$ observation density, with Laplacian noise. Results, averaged over $10$ random seeds, are shown in Figure~\ref{fig:synth}.

\begin{figure}[h]
\centering
\includegraphics[width=\textwidth]{all_tasks.pdf}
\caption{Convergence of the generalized Frank--Wolfe gap $\hat\Delta_k$ on the three synthetic tasks. Both proposed variants substantially outperform Vanilla SCFW and Clipped SCFW. The dashed line indicates the theoretical $\mathcal{O}(K^{-1/4})$ rate.}
\label{fig:synth}
\end{figure}

\textbf{Real-data experiments.}
Task~1 uses the \texttt{a9a} LIBSVM dataset \citep{chang2011libsvm}, partitioned into $m = 10$ disjoint subsets stratified by label, inducing significant task heterogeneity; we set $\tau = 10$. Task~2 uses daily closing-price returns of S\&P~500 constituents over 2014--2023 \citep{wrds_compustat}, restricted to the $d = 100$ most liquid assets continuously listed throughout the period; real returns are well known to exhibit heavy tails \citep{cont2001empirical}. Task~3 uses the \texttt{MovieLens-100K} rating matrix \citep{harper2015movielens}, where the user-rating distribution is heavy-tailed. Each experiment is averaged over $5$ random seeds. Results are shown in Figure~\ref{fig:real}.

\begin{figure}[h]
\centering
\includegraphics[width=\textwidth]{all_tasks_real.pdf}
\caption{Convergence on real-world datasets (\texttt{a9a}, S\&P~500, \texttt{MovieLens-100K}). The qualitative behavior tracks the synthetic study: Vanilla SCFW does not converge, Clipped SCFW reaches a tunable plateau, and both variants achieve the predicted $\mathcal{O}(K^{-1/4})$ rate.}
\label{fig:real}
\end{figure}

\textbf{Discussion.}
Across both synthetic (Figure~\ref{fig:synth}) and real-world (Figure~\ref{fig:real}) experiments, three observations are consistent. First, Vanilla SCFW fails to converge: the gap stagnates because plug-in stochastic estimates are biased through the non-linear $F$. Second, Clipped SCFW partially stabilizes the trajectory but plateaus well above the rate of the momentum-based variants and is sensitive to the clipping threshold. Third, both Variant I (Polyak) and Variant II (Taylor) achieve the predicted $\mathcal{O}(K^{-1/4})$ rate.

On Task~1, the two variants perform essentially identically; on Tasks~2 and~3, Variant I slightly outperforms Variant II in the regime explored. Both tasks satisfy Assumption~\ref{assump:inner_lipschitz}, so Variant I is applicable in either case. The schedules used here are the theoretically-prescribed ones; we did not perform per-task tuning of $\rho$, which may close the gap. Variant II remains the theoretically robust choice when Assumption~\ref{assump:inner_lipschitz} cannot be verified.

\section{Conclusion}
In this work, we introduced the first stochastic projection-free algorithm for fully composite optimization capable of handling non-smooth outer functions. Our approach achieves the minimax optimal $\mathcal{O}(K^{-1/4})$ rate for non-convex objectives and an $\mathcal{O}(K^{-1/3})$ rate for convex objectives under standard finite-variance noise assumptions. Furthermore, the algorithm seamlessly adapts to the noise level, yielding a clean reduction to optimal deterministic rates when the stochastic noise vanishes. 

\textbf{Limitations and Future Work.}
While our theoretical and empirical results demonstrate the efficacy of the proposed framework, several open challenges remain. First, our analysis assumes that the non-smooth outer function $F$ is deterministic; extending our framework to the \textit{doubly stochastic} regime—where $F$ is itself an expectation—remains an important hurdle. Second, our current convergence guarantees hold in expectation. Deriving high-probability bounds, particularly in settings corrupted by heavy-tailed noise, is a crucial next step for robust deployment. Finally, a highly compelling direction for future research is extending our framework to distributed and federated learning settings, where projection-free methods for non-smooth compositional objectives could significantly reduce communication bottlenecks.

\bibliographystyle{plainnat}
\bibliography{references}

\newpage
\appendix

\section{Auxiliary Inequalities} \label{sec:auxiliary}

\begin{lemma}[Norm power convexity] \label{lem:norm_convex}
For vectors $\va_1, \ldots, \va_m$ and $s \ge 1$: $\|\sum_i \va_i\|^s \le m^{s-1} \sum_i \|\va_i\|^s$.
\end{lemma}

\begin{lemma}[Subadditivity of the $r$-th root] \label{lem:subadd}
For $a_i \ge 0$ and $s \ge 1$: $(\sum_i a_i)^{1/s} \le \sum_i a_i^{1/s}$.
\end{lemma}

\begin{lemma}[Jensen for weighted sums] \label{lem:jensen}
If $w_i \ge 0$ with $\sum_i w_i \le 1$, then for any $a_i \ge 0$ and $s \ge 1$: $(\sum_i w_i a_i)^s \le \sum_i w_i a_i^s$.
\end{lemma}

\begin{lemma}[von Bahr--Esseen] \label{lem:vbe}
Let $\{\vX_i\}_{i=1}^k$ be a martingale difference sequence in $\R^n$. For $r \in (1, 2]$, there is a constant $C_r > 0$ such that $\E\|\sum_i \vX_i\|^r \le C_r \sum_i \E\|\vX_i\|^r$. At $r = 2$, $C_2 = 1$.
\end{lemma}

\section{Composite Curvature Bound} \label{sec:curvature_appendix}

\begin{proposition}[Bounded composite curvature] \label{prop:curvature}
Under Assumptions~\ref{assump:domain}--\ref{assump:inner_smooth}, the constant $\mathcal{S}$ defined in \eqref{eq:curvature_def} satisfies $\mathcal{S} \le L_F L D_{\cX}^2 \sqrt{n}$.
\end{proposition}

\begin{proof}
Fix $\vx, \vy \in \cX$ and $\gamma \in (0, 1]$, and set $\vy_\gamma = \vx + \gamma(\vy - \vx)$. By Assumption~\ref{assump:inner_smooth} applied componentwise, $|f_i(\vy_\gamma) - f_i(\vx) - \nabla f_i(\vx)^\top (\vy_\gamma - \vx)| \le \tfrac{L\gamma^2}{2} \|\vy - \vx\|^2$ for each $i$. Setting $\bm\Delta := \vf(\vy_\gamma) - \vf(\vx) - \nabla \vf(\vx)(\vy_\gamma - \vx)$, we get $\|\bm\Delta\| \le \tfrac{L\gamma^2 \sqrt{n}}{2} \|\vy - \vx\|^2$. By $L_F$-Lipschitzness of $F$ in $\vu$ and $\|\vy - \vx\| \le D_{\cX}$:
\[
F(\vf(\vy_\gamma), \vy_\gamma) - F(\vf(\vx) + \gamma\nabla\vf(\vx)(\vy - \vx), \vy_\gamma) \le L_F \|\bm\Delta\| \le \tfrac{L_F L \gamma^2 \sqrt{n} D_\cX^2}{2}.
\]
Multiplying by $2/\gamma^2$ gives the claim.
\end{proof}

\section{Full Statements: Progress and Tracking-Error Lemmas} \label{sec:lemmas_appendix}

Let the tracking errors be $\bm{\delta}_{g,k} := \mV_k - \nabla \vf(\vy_k)$ and $\bm{\delta}_{f,k} := \vz_k - \vf(\vy_k)$.

\begin{lemma}[GLMO progress with tracking errors] \label{lem:glmo_main}
Under Assumptions~\ref{assump:domain}, \ref{assump:outer}, \ref{assump:inner_smooth}:
\[
\varphi(\vy_{k+1}) \le \varphi(\vy_k) - \gamma_k \hat\Delta_k + \tfrac{\gamma_k^2}{2} \mathcal{S} + 2 \gamma_k L_F (\|\bm{\delta}_{f,k}\| + D_{\cX} \|\bm{\delta}_{g,k}\|).
\]
\end{lemma}

\begin{lemma}[Jacobian tracker, $r$-th moment] \label{lem:grad_tracker_main}
Under Assumptions~\ref{assump:inner_smooth} and \ref{assump:noise}, with constant $\beta \in (0,1]$ and $\|\vy_k - \vy_{k-1}\| \le \gamma D_{\cX}$:
\[
\E\|\bm{\delta}_{g,k}\|^r \le 3^{r-1}\!\left[(1-\beta)^{rk} \E\|\bm{\delta}_{g,0}\|^r + \frac{L^r D_{\cX}^r \gamma^r}{\beta^r} + C_r \sigma_g^r \beta^{r-1}\right].
\]
\end{lemma}

\begin{lemma}[Polyak function tracker, $r$-th moment] \label{lem:func_polyak_main}
Under Assumptions~\ref{assump:inner_smooth}, \ref{assump:noise}, \ref{assump:inner_lipschitz}, with constant $\rho \in (0,1]$ and $\|\vy_k - \vy_{k-1}\| \le \gamma D_{\cX}$:
\[
\E\|\bm{\delta}_{f,k}\|^r \le 3^{r-1}\!\left[(1-\rho)^{rk} \E\|\bm{\delta}_{f,0}\|^r + \frac{G^r D_{\cX}^r \gamma^r}{\rho^r} + C_r \sigma_f^r \rho^{r-1}\right].
\]
\end{lemma}

\begin{lemma}[Taylor function tracker, $r$-th moment] \label{lem:func_taylor_main}
Under Assumptions~\ref{assump:inner_smooth} and \ref{assump:noise}, with constant $\rho \in (0,1]$ and $\|\vy_k - \vy_{k-1}\| \le \gamma D_{\cX}$:
\begin{align*}
\E\|\bm{\delta}_{f,k}\|^r &\le 4^{r-1}\!\Big[(1-\rho)^{rk} \E\|\bm{\delta}_{f,0}\|^r + \frac{L^r D_{\cX}^{2r} \gamma^{2r}}{2^r \rho^r} \\
&\qquad + \frac{D_{\cX}^r \gamma^r}{\rho^r} \max_{1 \le i \le k} \E\|\bm{\delta}_{g,i}\|^r + C_r \sigma_f^r \rho^{r-1}\Big].
\end{align*}
\end{lemma}

The Taylor correction reduces the dependence of the function-tracker bias on the step size from $\gamma^r/\rho^r$ (Polyak) to $\gamma^{2r}/\rho^r$ (Taylor), at the cost of a coupling to the Jacobian-tracker error which is itself controlled.

\section{Proof of Lemma~\ref{lem:glmo_main}} \label{sec:proof_glmo}

\begin{proof}
Define the exact and tracked surrogates $l_k(\vx) := \vf(\vy_k) + \nabla \vf(\vy_k)(\vx - \vy_k)$ and $\tilde l_k(\vx) := \vz_k + \mV_k(\vx - \vy_k)$. Their difference at any $\vx \in \cX$ is uniformly bounded:
\begin{equation} \label{eq:uniform_surrogate_bound}
\|\tilde l_k(\vx) - l_k(\vx)\| \le \|\bm{\delta}_{f,k}\| + D_{\cX} \|\bm{\delta}_{g,k}\| =: E_k.
\end{equation}

\textbf{Step 1.} By definition of $\mathcal{S}$, applied with $\vx \leftarrow \vy_k$, $\vy \leftarrow \vx_{k+1}$, $\gamma \leftarrow \gamma_k$:
\begin{equation} \label{eq:step1_curvature}
\varphi(\vy_{k+1}) \le F(\vf(\vy_k) + \gamma_k \nabla \vf(\vy_k)(\vx_{k+1} - \vy_k), \vy_{k+1}) + \tfrac{\gamma_k^2}{2} \mathcal{S}.
\end{equation}

\textbf{Step 2.} Using $\vf(\vy_k) + \gamma_k \nabla \vf(\vy_k)(\vx_{k+1} - \vy_k) = (1-\gamma_k)\vf(\vy_k) + \gamma_k l_k(\vx_{k+1})$, joint subhomogeneity, and monotonicity:
\begin{equation} \label{eq:step2_subhomog}
F(\cdot, \vy_{k+1}) \le (1 - \gamma_k) \varphi(\vy_k) + \gamma_k F(l_k(\vx_{k+1}), \vx_{k+1}).
\end{equation}

\textbf{Step 3.} Let $\bar \vx_{k+1} \in \argmin_{\vx \in \cX} F(l_k(\vx), \vx)$, so $F(l_k(\bar\vx_{k+1}), \bar\vx_{k+1}) = \varphi(\vy_k) - \hat\Delta_k$. By two applications of $L_F$-Lipschitzness and the GLMO optimality of $\vx_{k+1}$:
\begin{align*}
F(l_k(\vx_{k+1}), \vx_{k+1})
&\le F(\tilde l_k(\vx_{k+1}), \vx_{k+1}) + L_F E_k
\le F(\tilde l_k(\bar \vx_{k+1}), \bar \vx_{k+1}) + L_F E_k \\
&\le F(l_k(\bar \vx_{k+1}), \bar \vx_{k+1}) + 2 L_F E_k = \varphi(\vy_k) - \hat\Delta_k + 2 L_F E_k.
\end{align*}

\textbf{Step 4.} Substituting into \eqref{eq:step2_subhomog} and \eqref{eq:step1_curvature}:
\[
\varphi(\vy_{k+1}) \le \varphi(\vy_k) - \gamma_k \hat\Delta_k + \tfrac{\gamma_k^2}{2} \mathcal{S} + 2 \gamma_k L_F E_k. \qedhere
\]
\end{proof}

\section{Proof of Lemma~\ref{lem:grad_tracker_main}} \label{sec:proof_grad}

\begin{proof}
Let $\bm{\xi}_{g,i} := \nabla \tilde \vf(\vy_i; \xi_i) - \nabla \vf(\vy_i)$ and $\bm{\Delta}_i := \nabla \vf(\vy_{i-1}) - \nabla \vf(\vy_i)$. By Assumptions~\ref{assump:noise}, \ref{assump:inner_smooth}: $\E[\bm\xi_{g,i} | \mathcal{F}_{i-1}] = \mathbf{0}$, $\E\|\bm\xi_{g,i}\|^r \le \sigma_g^r$, $\|\bm\Delta_i\| \le L\gamma D_\cX$.

\emph{Step 1: Recursion.} Subtracting $\nabla \vf(\vy_k)$ from both sides of the Polyak update,
\[
\bm{\delta}_{g,k} = (1-\beta)\bm{\delta}_{g,k-1} + (1-\beta) \bm{\Delta}_k + \beta \bm{\xi}_{g,k}.
\]

\emph{Step 2: Unrolling.} Iterating from $i = 1$ to $k$:
\[
\bm{\delta}_{g,k} = \underbrace{(1-\beta)^k \bm{\delta}_{g,0}}_{T_1} + \underbrace{\sum_{i=1}^k (1-\beta)^{k-i+1} \bm\Delta_i}_{T_2} + \underbrace{\beta \sum_{i=1}^k (1-\beta)^{k-i} \bm\xi_{g,i}}_{T_3}.
\]

\emph{Step 3:} By Lemma~\ref{lem:norm_convex}, $\|\bm\delta_{g,k}\|^r \le 3^{r-1}(\|T_1\|^r + \|T_2\|^r + \|T_3\|^r)$.

\emph{Step 4: $T_1$.} $\|T_1\|^r = (1-\beta)^{rk} \|\bm\delta_{g,0}\|^r$.

\emph{Step 5: $T_2$.} By the triangle inequality and the geometric series, $\|T_2\| \le L\gamma D_\cX/\beta$, so $\|T_2\|^r \le L^r\gamma^r D_\cX^r/\beta^r$.

\emph{Step 6: $T_3$.} The sequence $\{(1-\beta)^{k-i}\bm\xi_{g,i}\}$ is a martingale difference sequence. By Lemma~\ref{lem:vbe} and using $(1-\beta)^r \le 1-\beta$:
\[
\E\|T_3\|^r \le \beta^r C_r \sum_i (1-\beta)^{r(k-i)} \sigma_g^r \le \beta^r C_r \sigma_g^r / \beta = C_r \sigma_g^r \beta^{r-1}.
\]

\emph{Step 7:} Combining, $\E\|\bm\delta_{g,k}\|^r \le 3^{r-1}\big[(1-\beta)^{rk}\E\|\bm\delta_{g,0}\|^r + L^r\gamma^r D_\cX^r/\beta^r + C_r \sigma_g^r \beta^{r-1}\big]$.
\end{proof}

\section{Proof of Lemma~\ref{lem:func_polyak_main}} \label{sec:proof_polyak}

\begin{proof}
Structurally identical to the proof of Lemma~\ref{lem:grad_tracker_main}, replacing $\beta \leftarrow \rho$, $\nabla \tilde\vf \leftarrow \tilde\vf$, $\nabla \vf \leftarrow \vf$, and using Assumption~\ref{assump:inner_lipschitz} for the deterministic drift: $\bm\Delta_i = \vf(\vy_{i-1}) - \vf(\vy_i)$ satisfies $\|\bm\Delta_i\| \le G\gamma D_\cX$.
\end{proof}

\section{Proof of Lemma~\ref{lem:func_taylor_main}} \label{sec:proof_taylor}

\begin{proof}
\emph{Step 1: Algebraic identity.} Define the Taylor remainder $\vR_k := \vf(\vy_k) - \vf(\vy_{k-1}) - \nabla\vf(\vy_k)(\vy_k - \vy_{k-1})$, with $\|\vR_k\| \le L\gamma^2 D_\cX^2/2$ by Assumption~\ref{assump:inner_smooth}. Subtracting $\vf(\vy_k)$ from the Variant II update and substituting $\vf(\vy_k) = \vf(\vy_{k-1}) + \nabla\vf(\vy_k)(\vy_k - \vy_{k-1}) - \vR_k$ inside the bracket:
\begin{equation}\label{eq:taylor_recursion_appendix}
\bm{\delta}_{f,k} = (1-\rho) \bm{\delta}_{f,k-1} + (1-\rho) \bm{\delta}_{g,k}(\vy_k - \vy_{k-1}) + (1-\rho) \vR_k + \rho \bm{\xi}_{f,k}.
\end{equation}

\emph{Step 2: Unrolling.} Iterating \eqref{eq:taylor_recursion_appendix} from $i = 1$ to $k$:
\begin{align*}
\bm\delta_{f,k} &= \underbrace{(1-\rho)^k \bm{\delta}_{f,0}}_{T_1} + \underbrace{\sum_{i=1}^k (1-\rho)^{k-i+1} \bm{\delta}_{g,i}(\vy_i - \vy_{i-1})}_{T_2} \\
&\quad + \underbrace{\sum_{i=1}^k (1-\rho)^{k-i+1} \vR_i}_{T_3} + \underbrace{\rho \sum_{i=1}^k (1-\rho)^{k-i} \bm\xi_{f,i}}_{T_4}.
\end{align*}

\emph{Step 3:} By Lemma~\ref{lem:norm_convex}, $\E\|\bm\delta_{f,k}\|^r \le 4^{r-1}(\E\|T_1\|^r + \E\|T_2\|^r + \E\|T_3\|^r + \E\|T_4\|^r)$.

\emph{Step 4: $T_1$.} $\E\|T_1\|^r = (1-\rho)^{rk}\E\|\bm\delta_{f,0}\|^r$.

\emph{Step 5: $T_3$ (Taylor remainder drift).} $\|T_3\| \le L\gamma^2 D_\cX^2/(2\rho)$, hence $\E\|T_3\|^r \le L^r\gamma^{2r} D_\cX^{2r}/(2^r\rho^r)$.

\emph{Step 6: $T_4$ (martingale noise).} As in the Jacobian-tracker proof, $\E\|T_4\|^r \le C_r \sigma_f^r \rho^{r-1}$.

\emph{Step 7: $T_2$ (cross-tracking term).} Define weights $w_i := \rho(1-\rho)^{k-i}$, so $\sum_i w_i \le 1$. Then $T_2 = \tfrac{1-\rho}{\rho}\sum_i w_i \bm\delta_{g,i}(\vy_i - \vy_{i-1})$, and $\|T_2\| \le \tfrac{\gamma D_\cX}{\rho}\sum_i w_i \|\bm\delta_{g,i}\|$. By Lemma~\ref{lem:jensen}:
\[
\E\|T_2\|^r \le \tfrac{\gamma^r D_\cX^r}{\rho^r} \sum_i w_i \E\|\bm\delta_{g,i}\|^r \le \tfrac{\gamma^r D_\cX^r}{\rho^r} \max_{1 \le i \le k} \E\|\bm\delta_{g,i}\|^r.
\]

\emph{Step 8:} Combining yields the lemma.
\end{proof}

\section{Proof of Theorem~\ref{thm:nonconvex}} \label{sec:proof_nonconvex}

We give a unified argument that handles both variants. Throughout, $C_r$ denotes the von Bahr--Esseen constant.

\subsection{First-moment tracking-error bounds}

For any non-negative random variable $X$, Jensen gives $\E[X] \le (\E X^r)^{1/r}$; combined with Lemma~\ref{lem:subadd}, the $r$-th-moment bounds of Lemmas~\ref{lem:grad_tracker_main}--\ref{lem:func_taylor_main} imply, with the schedule $\beta = \rho$:
\begin{align}
\E\|\bm\delta_{g,k}\| &\le \mathcal{E}_{g,0}(1-\beta)^k + U_g \cdot \tfrac{\gamma}{\beta} + U_g \cdot \beta^{(r-1)/r}, \label{eq:grad_first_simplified} \\
\E\|\bm\delta_{f,k}\|_{\mathrm{Var\ I}} &\le \mathcal{E}_{f,0}(1-\rho)^k + U_f^{\mathrm{I}} \cdot \tfrac{\gamma}{\rho} + U_f^{\mathrm{I}} \cdot \rho^{(r-1)/r}, \label{eq:func_polyak_first} \\
\E\|\bm\delta_{f,k}\|_{\mathrm{Var\ II}} &\le \mathcal{E}_{f,0}^{\mathrm{II}}(1-\rho)^k + U_f^{\mathrm{II}}\Big(\tfrac{\gamma^2}{\rho} + \tfrac{\gamma}{\rho}\mathcal{E}_{g,0} + \tfrac{\gamma^2}{\rho^2} + \tfrac{\gamma}{\rho^{1/r}} + \rho^{(r-1)/r}\Big), \label{eq:func_taylor_first}
\end{align}
where
\begin{align*}
\mathcal{E}_{g,0} &:= 3^{(r-1)/r}(\E\|\bm\delta_{g,0}\|^r)^{1/r}, &
U_g &:= 3^{(r-1)/r}\max\{LD_\cX, C_r^{1/r}\sigma_g\}, \\
\mathcal{E}_{f,0} &:= 3^{(r-1)/r}(\E\|\bm\delta_{f,0}\|^r)^{1/r}, &
U_f^{\mathrm{I}} &:= 3^{(r-1)/r}\max\{GD_\cX, C_r^{1/r}\sigma_f\}, \\
\mathcal{E}_{f,0}^{\mathrm{II}} &:= 4^{(r-1)/r}(\E\|\bm\delta_{f,0}\|^r)^{1/r}, &
U_f^{\mathrm{II}} &:= 4^{(r-1)/r}\max\{LD_\cX^2/2,\ U_g D_\cX,\ C_r^{1/r}\sigma_f\}.
\end{align*}

\subsection{Telescoping and balancing}

Summing Lemma~\ref{lem:glmo_main} from $k=0$ to $K-1$, telescoping using $\E[\varphi(\vy_K)] - \varphi(\vy_0) \ge -\Phi_0$, and dividing by $\gamma K$:
\begin{equation} \label{eq:telescope_main}
\min_{0 \le k \le K-1} \E[\hat\Delta_k] \le \frac{\Phi_0}{\gamma K} + \frac{\gamma\mathcal{S}}{2} + \frac{2L_F}{K}\sum_{k=0}^{K-1}\big(\E\|\bm\delta_{f,k}\| + D_\cX \E\|\bm\delta_{g,k}\|\big).
\end{equation}

Each tracker bound has the form $A(1-\beta)^k + (\text{constant terms})$, and $\sum_{k=0}^{K-1}(1-\beta)^k \le 1/\beta$. With the schedule $\gamma = K^{-(2r-1)/(3r-2)}, \beta = \rho = K^{-r/(3r-2)}$, every term in \eqref{eq:telescope_main} can be verified to scale at most as $K^{-(r-1)/(3r-2)}$:
\begin{itemize}
\item $\Phi_0/(\gamma K) = K^{-(r-1)/(3r-2)}$ (tight); $\gamma \le K^{-(r-1)/(3r-2)}$ for $r \ge 1$;
\item $1/(K\beta) = K^{-(2r-2)/(3r-2)}$ (faster);
\item $\gamma/\beta = K^{-(r-1)/(3r-2)}$ (tight); $\beta^{(r-1)/r} = K^{-(r-1)/(3r-2)}$ (tight);
\item Variant II auxiliary terms $\gamma^2/\rho = K^{-1}$, $\gamma^2/\rho^2 = K^{-2(r-1)/(3r-2)}$, $\gamma/\rho^{1/r} = K^{-(2r-2)/(3r-2)}$ (all faster).
\end{itemize}

\subsection{Explicit constants}

Collecting prefactors:
\begin{equation} \label{eq:MI_explicit}
\boxed{M_{\mathrm{I}} = \Phi_0 + \tfrac{1}{2} L_F L D_{\cX}^2 \sqrt{n} + 2 L_F (\mathcal{E}_{f,0} + 2 U_f^{\mathrm{I}}) + 2 L_F D_{\cX}(\mathcal{E}_{g,0} + 2 U_g)}
\end{equation}
\begin{equation} \label{eq:MII_explicit}
\boxed{M_{\mathrm{II}} = \Phi_0 + \tfrac{1}{2} L_F L D_{\cX}^2 \sqrt{n} + 2 L_F (\mathcal{E}_{f,0}^{\mathrm{II}} + (3 + \mathcal{E}_{g,0}) U_f^{\mathrm{II}}) + 2 L_F D_{\cX}(\mathcal{E}_{g,0} + 2 U_g)}
\end{equation}

For each variant, $\min_k \E[\hat\Delta_k] \le M_{(\cdot)} \cdot K^{-(r-1)/(3r-2)}$. This completes the proof.\hfill\qed

\section{Proof of Theorem~\ref{thm:convex}} \label{sec:proof_convex}

We work with primal suboptimality $h_k := \E[\varphi(\vy_k) - \varphi(\vx^*)]$.

\subsection{Primal gap inequality}

\begin{lemma}[Primal suboptimality from FW gap, convex case] \label{lem:convex_primal}
Under Assumptions~\ref{assump:outer} and \ref{assump:convex}: $\hat\Delta(\vy) \ge \varphi(\vy) - \varphi(\vx^*)$.
\end{lemma}

\begin{proof}
Let $\bar\vx \in \argmin_{\vx \in \cX} F(l(\vx), \vx)$ where $l(\vx) := \vf(\vy) + \nabla\vf(\vy)(\vx - \vy)$. Since $\bar\vx$ is the minimizer over $\cX$,
$F(l(\bar\vx), \bar\vx) \le F(l(\vx^*), \vx^*)$. By convexity of each $f_i$ (Assumption~\ref{assump:convex}), $\vf(\vx^*) \ge l(\vx^*)$ componentwise. By monotonicity of $F$ in $\vu$, $F(l(\vx^*), \vx^*) \le F(\vf(\vx^*), \vx^*) = \varphi(\vx^*)$. Hence $\hat\Delta(\vy) \ge \varphi(\vy) - \varphi(\vx^*)$.
\end{proof}

Combining with Lemma~\ref{lem:glmo_main}:
\begin{equation} \label{eq:convex_recur}
h_{k+1} \le (1 - \gamma_k) h_k + \tfrac{\gamma_k^2}{2}\mathcal{S} + 2\gamma_k L_F\big(\E\|\bm\delta_{f,k}\| + D_\cX \E\|\bm\delta_{g,k}\|\big).
\end{equation}

\subsection{Tracking errors under decreasing schedule}

\begin{lemma}[Decreasing-schedule tracking] \label{lem:convex_tracking}
Set $y := r/(2r-1)$ and $q := y(r-1) = r(r-1)/(2r-1)$. Choose $c_0 = 1$ and $k_0 := \lceil (4q)^{1/(1-y)}\rceil + 2$. Under Assumptions~\ref{assump:inner_smooth}, \ref{assump:noise} and the schedule \eqref{eq:convex_schedule}, there exist explicit constants $Q_g, Q_f > 0$ such that for all $k \ge 0$:
\[
\E\|\bm\delta_{g,k}\| \le Q_g (k+k_0)^{-(r-1)/(2r-1)}, \qquad \E\|\bm\delta_{f,k}\| \le Q_f (k+k_0)^{-(r-1)/(2r-1)}.
\]
\end{lemma}

\begin{proof}
We prove the Jacobian bound; the function tracker is analogous. From the recursion of Lemma~\ref{lem:grad_tracker_main}, applying the Young-type inequality $\|\va + \vb\|^r \le (1+\eta)^{r-1}\|\va\|^r + (1+1/\eta)^{r-1}\|\vb\|^r$ with $\eta = \beta_k$ and using $(1-\beta_k)^r(1+\beta_k)^{r-1} \le 1-\beta_k$:
\[
\E\|\bm\delta_{g,k}\|^r \le (1-\beta_k)\E\|\bm\delta_{g,k-1}\|^r + 2^{r-1}\beta_k^{-(r-1)} L^r D_\cX^r \gamma_{k-1}^r + C_r \sigma_g^r \beta_k^r.
\]
With $\gamma_{k-1} \le 4(k+k_0)^{-1}$ and $\beta_k = (k+k_0)^{-y}$, both $\beta_k^{-(r-1)}\gamma_{k-1}^r$ and $\beta_k^r$ scale as $(k+k_0)^{-yr}$ (since $y(r-1) - r = -yr$):
\begin{equation}\label{eq:convex_recur_clean}
\E\|\bm\delta_{g,k}\|^r \le \big(1 - (k+k_0)^{-y}\big)\E\|\bm\delta_{g,k-1}\|^r + M_g (k+k_0)^{-yr},
\end{equation}
where $M_g := 2^{r-1}\cdot 4^r L^r D_\cX^r + C_r \sigma_g^r$.

\emph{Induction.} We prove $\E\|\bm\delta_{g,k}\|^r \le \bar Q_g(k+k_0)^{-q}$ by induction. Base case: $\bar Q_g \ge k_0^q \E\|\bm\delta_{g,0}\|^r$. For the inductive step, using $(1-x)^{-q} \le 1 + 2qx$ for $x \le 1/2$ and $q \le 1$ (which holds for $k_0 \ge 2$):
\[
(k-1+k_0)^{-q} \le (k+k_0)^{-q}\big(1 + 2q/(k+k_0)\big).
\]
Substituting and using $-q-y = -yr$:
\[
\E\|\bm\delta_{g,k}\|^r \le \bar Q_g(k+k_0)^{-q} + (k+k_0)^{-yr}[M_g - \bar Q_g] + 2q\bar Q_g(k+k_0)^{-q-1}.
\]
Since $(k+k_0)^{-q-1} \le k_0^{y-1}(k+k_0)^{-yr}$, the induction closes provided $\bar Q_g \ge M_g/(1 - 2q k_0^{y-1})$. Choosing $k_0$ as in the lemma ensures $2qk_0^{y-1} \le 1/2$, so $\bar Q_g \ge 2 M_g$ suffices:
\begin{equation} \label{eq:Qg_explicit}
\bar Q_g := \max\{k_0^q \E\|\bm\delta_{g,0}\|^r,\ 2 M_g\}, \qquad Q_g := \bar Q_g^{1/r}.
\end{equation}

The function tracker satisfies an analogous recurrence with the additional cross-coupling term involving $\E\|\bm\delta_{g,k}\|^r$; substituting the just-proved Jacobian bound, the same induction yields
\begin{equation} \label{eq:Qf_explicit}
Q_f := \big(\max\{k_0^q \E\|\bm\delta_{f,0}\|^r,\ 2 M_f\}\big)^{1/r},
\end{equation}
where $M_f := 4^{r-1}[4^r L^r D_\cX^{2r}/2^r + 4^r D_\cX^r \bar Q_g + C_r \sigma_f^r]$ (Variant II); for Variant I, replace $L^r D_\cX^{2r}/2^r + D_\cX^r \bar Q_g$ with $G^r D_\cX^r$. Applying Jensen yields the first-moment bounds.
\end{proof}

\subsection{Solving the primal recurrence}

Let $p := (r-1)/(2r-1)$. Substituting Lemma~\ref{lem:convex_tracking} into \eqref{eq:convex_recur} and using $k_0 \ge 2$ to bound $(k+k_0)^{-p} \le 2^p (k+2)^{-p}$:
\[
h_{k+1} \le \big(1 - \tfrac{2}{k+2}\big) h_k + \frac{B}{(k+2)^{1+p}}, \qquad B := 2\mathcal{S} + 4 L_F (Q_f + D_\cX Q_g)\cdot 2^p.
\]

\emph{Induction claim:} $h_k \le A_{\mathrm{cvx}}(k+2)^{-p}$. Base case: $A_{\mathrm{cvx}} \ge 2^p \Phi_0$. For the inductive step, using $(k+3)^{-p} \ge (k+2)^{-p}(1 - p/(k+2))$:
\[
h_{k+1} \le A_{\mathrm{cvx}}(k+2)^{-p} - \frac{2A_{\mathrm{cvx}} - B}{(k+2)^{1+p}} \le A_{\mathrm{cvx}}(k+3)^{-p}
\]
provided $A_{\mathrm{cvx}}(2-p) \ge B$. Taking
\begin{equation}\label{eq:Acvx_def}
\boxed{A_{\mathrm{cvx}} := \max\Big\{2^p \Phi_0,\ \tfrac{B}{2-p}\Big\}}
\end{equation}
closes the induction. Since $\hat\Delta_K \ge h_K$ by Lemma~\ref{lem:convex_primal}, $\E[\hat\Delta_K] \le A_{\mathrm{cvx}}(K+2)^{-p}$.\hfill\qed

\section{Proof of Corollary~\ref{cor:deterministic}} \label{sec:proof_deterministic}

\begin{proof}
With $\sigma_f = \sigma_g = 0$ and $\beta_k = \rho_k = 1$, the trackers $\mV_k = \nabla\tilde\vf(\vy_k; \xi_k) = \nabla\vf(\vy_k)$ (for Variant I); for Variant II, the Taylor correction term has prefactor $1 - \rho_k = 0$, and $\vz_k = \tilde\vf(\vy_k; \xi_k) = \vf(\vy_k)$. Hence $\bm\delta_{g,k} = \bm\delta_{f,k} = \mathbf{0}$ for all $k \ge 1$. Lemma~\ref{lem:glmo_main} reduces to
\[
\varphi(\vy_{k+1}) \le \varphi(\vy_k) - \gamma_k \hat\Delta_k + \tfrac{\gamma_k^2}{2}\mathcal{S},
\]
which is exactly the deterministic progress bound (Lemma~A.1 of \citet{vladarean2023first}). Their Theorems~3.1 and 3.2 then apply with the cited deterministic schedules.
\end{proof}

\section{Acceleration via STORM} \label{sec:storm_extension}

The Polyak Jacobian tracker yields rate $\mathcal{O}(K^{-(r-1)/(3r-2)})$. STORM replaces the Polyak update by a stochastic finite-difference correction that achieves $\mathcal{O}(K^{-(r-1)/(2r-1)})$, at the cost of $r$-average smoothness.

\begin{assumption}[$r$-average smoothness] \label{assump:avg_smooth}
There exists $\bar L > 0$ such that $\E_\xi\|\nabla\tilde\vf(\vx; \xi) - \nabla\tilde\vf(\vy; \xi)\|^r \le \bar L^r \|\vx - \vy\|^r$ for all $\vx, \vy \in \cX$.
\end{assumption}

The STORM Jacobian update reuses $\xi_k$ at $\vy_k$ and $\vy_{k-1}$:
\begin{equation} \label{eq:storm_grad}
\mV_k = \nabla\tilde\vf(\vy_k; \xi_k) + (1-\beta)[\mV_{k-1} - \nabla\tilde\vf(\vy_{k-1}; \xi_k)].
\end{equation}
The function tracker remains the Taylor-corrected Variant II.

\begin{lemma}[STORM Jacobian tracking error] \label{lem:storm_tracker}
Under Assumptions~\ref{assump:inner_smooth}, \ref{assump:noise}, \ref{assump:avg_smooth} with \eqref{eq:storm_grad} and $\|\vy_k - \vy_{k-1}\| \le \gamma D_\cX$:
\[
\E\|\bm\delta_{g,k}\|^r \le (1-\beta)^k \E\|\bm\delta_{g,0}\|^r + \frac{2^r C_r \bar L^r D_\cX^r \gamma^r}{\beta} + C_r \sigma_g^r \beta^{r-1}.
\]
\end{lemma}

\begin{proof}
Subtracting $\nabla\vf(\vy_k)$ from \eqref{eq:storm_grad}: $\bm\delta_{g,k} = (1-\beta)\bm\delta_{g,k-1} + (1-\beta)\bm\zeta_k + \beta\bm\xi_{g,k}$, where $\bm\zeta_k := [\nabla\tilde\vf(\vy_k; \xi_k) - \nabla\tilde\vf(\vy_{k-1}; \xi_k)] - [\nabla\vf(\vy_k) - \nabla\vf(\vy_{k-1})]$. Both $\bm\zeta_k$ and $\bm\xi_{g,k}$ are mean-zero conditional on $\mathcal{F}_{k-1}$. By Assumption~\ref{assump:avg_smooth}, $\E\|\bm\zeta_k\|^r \le 2^r \bar L^r \gamma^r D_\cX^r$. Conditioning and applying Lemma~\ref{lem:vbe} with $(1-\beta)^r \le 1-\beta$, then iterating, yields the claim.
\end{proof}

\begin{theorem}[STORM acceleration] \label{thm:storm_rate}
Under Assumptions~\ref{assump:domain}--\ref{assump:noise} and \ref{assump:avg_smooth}, with the STORM update \eqref{eq:storm_grad}, Variant II function tracker, and constant schedule $\gamma = \beta = \rho = K^{-r/(2r-1)}$:
\[
\min_{1 \le k \le K} \E[\hat\Delta_k] = \mathcal{O}\!\left(K^{-(r-1)/(2r-1)}\right).
\]
At $r = 2$, the rate is $\mathcal{O}(K^{-1/3})$.
\end{theorem}

\begin{proof}
By Jensen on Lemma~\ref{lem:storm_tracker}: the drift $(\gamma^r/\beta)^{1/r} = \gamma\beta^{-1/r}$ and noise $\beta^{(r-1)/r}$ both scale as $K^{-(r-1)/(2r-1)}$ under the schedule. The Taylor function tracker (Lemma~\ref{lem:func_taylor_main}) inherits this rate via its cross-coupling term. Substituting into the GLMO progress bound and dividing by $\gamma K$, the initialization gap $1/(\gamma K) = K^{-(r-1)/(2r-1)}$ matches; the curvature term $\gamma$ decays faster.
\end{proof}

\section{Hessian-Corrected Momentum and the Average-Smoothness Equivalence} \label{sec:hessian_extension}

We now present an alternative acceleration mechanism via a randomized Hessian--vector correction, and prove that bounded Hessian-noise moments imply $r$-average smoothness with explicit constants.

\subsection{Bounded Hessian Noise Implies Average Smoothness}

\begin{assumption}[Bounded $r$-th moment of stochastic Hessian] \label{assump:hessian_var}
There exist $\sigma_H \ge 0$ and $r \in (1,2]$ such that $\tilde\vf(\cdot; \xi)$ is twice differentiable on $\cX$ a.s., with $\E_\xi\|\nabla^2\tilde\vf(\vx; \xi) - \nabla^2\vf(\vx)\|_{\mathrm{op}}^r \le \sigma_H^r$ for all $\vx \in \cX$.
\end{assumption}

\begin{proposition}[Hessian noise implies $r$-average smoothness] \label{prop:hessian_implies_avg}
Under Assumption~\ref{assump:inner_smooth} (which implies $\|\nabla^2\vf(\vx)\|_{\mathrm{op}} \le L$ when $\vf$ is $C^2$) and Assumption~\ref{assump:hessian_var}, Assumption~\ref{assump:avg_smooth} holds with $\bar L^r = 2^{r-1}(\sigma_H^r + L^r)$.
\end{proposition}

\begin{proof}
\emph{Step 1.} Using $\|\vA\|^r \le 2^{r-1}(\|\vA - \vB\|^r + \|\vB\|^r)$:
\[
\E_\xi\|\nabla^2\tilde\vf(\vx; \xi)\|_{\mathrm{op}}^r \le 2^{r-1}\sigma_H^r + 2^{r-1} L^r =: \bar L^r.
\]

\emph{Step 2.} The mean-value theorem gives, for each $\xi$ a.s.,
\[
\nabla\tilde\vf(\vx;\xi) - \nabla\tilde\vf(\vy;\xi) = \int_0^1 \nabla^2\tilde\vf(\vy + t(\vx - \vy); \xi)[\vx - \vy]\,dt.
\]

\emph{Step 3.} Taking the $r$-th power and applying Jensen on the uniform measure on $[0,1]$ (since $s \mapsto s^r$ is convex):
\[
\|\nabla\tilde\vf(\vx;\xi) - \nabla\tilde\vf(\vy;\xi)\|^r \le \int_0^1 \|\nabla^2\tilde\vf(\vy + t(\vx - \vy); \xi)\|_{\mathrm{op}}^r\,dt \cdot \|\vx - \vy\|^r.
\]

\emph{Step 4.} Taking expectation and applying Fubini:
\[
\E_\xi\|\nabla\tilde\vf(\vx; \xi) - \nabla\tilde\vf(\vy; \xi)\|^r \le \int_0^1 \E_\xi\|\nabla^2\tilde\vf(\cdot; \xi)\|_{\mathrm{op}}^r\,dt \cdot \|\vx - \vy\|^r \le \bar L^r \|\vx - \vy\|^r.\qedhere
\]
\end{proof}

\subsection{Hessian-Corrected Algorithm and Convergence}

Let $\alpha_k \sim \mathrm{U}(0,1)$ and $\zeta_k$ be drawn independently of $\xi_k$ from the Hessian oracle. Define $\vy_{\alpha,k} := \vy_{k-1} + \alpha_k(\vy_k - \vy_{k-1})$ and $\tilde\mH_k := \nabla^2\tilde\vf(\vy_{\alpha,k}; \zeta_k)$. The Hessian-corrected Jacobian update is
\begin{equation} \label{eq:hessian_grad}
\mV_k = (1-\beta)\mV_{k-1} + \beta\nabla\tilde\vf(\vy_k; \xi_k) + (1-\beta)\tilde\mH_k(\vy_k - \vy_{k-1}).
\end{equation}

\begin{lemma}[Unbiasedness of Hessian--vector correction] \label{lem:hessian_unbiased}
Conditional on $\mathcal{F}_{k-1}$: $\E[\tilde\mH_k(\vy_k - \vy_{k-1}) | \mathcal{F}_{k-1}] = \nabla\vf(\vy_k) - \nabla\vf(\vy_{k-1})$.
\end{lemma}

\begin{proof}
By the fundamental theorem of calculus, $\nabla\vf(\vy_k) - \nabla\vf(\vy_{k-1}) = \E_\alpha[\nabla^2\vf(\vy_{\alpha,k})](\vy_k - \vy_{k-1})$. By unbiasedness of the stochastic Hessian and conditional independence of $\alpha_k, \zeta_k$:
\[
\E[\tilde\mH_k | \mathcal{F}_{k-1}, \alpha_k] = \nabla^2\vf(\vy_{\alpha,k}),
\]
and taking iterated expectation completes the proof.
\end{proof}

\begin{lemma}[Hessian-corrected Jacobian tracking error] \label{lem:hessian_track}
Under Assumptions~\ref{assump:inner_smooth}, \ref{assump:noise}, \ref{assump:hessian_var} with $\xi_k \perp \zeta_k | \mathcal{F}_{k-1}$ and $\|\vy_k - \vy_{k-1}\| \le \gamma D_\cX$:
\[
\E\|\bm\delta_{g,k}\|^r \le (1-\beta)^k \E\|\bm\delta_{g,0}\|^r + \frac{2^{r-1} C_r \bar L^r D_\cX^r \gamma^r}{\beta} + C_r \sigma_g^r \beta^{r-1},
\]
where $\bar L^r = 2^{r-1}(\sigma_H^r + L^r)$.
\end{lemma}

\begin{proof}
\emph{Step 1: Decomposition.} Subtracting $\nabla\vf(\vy_k)$ from \eqref{eq:hessian_grad}:
\[
\bm\delta_{g,k} = (1-\beta)\bm\delta_{g,k-1} + (1-\beta)\bm\zeta_{H,k} + \beta\bm\xi_{g,k},
\]
where $\bm\zeta_{H,k} := \tilde\mH_k(\vy_k - \vy_{k-1}) - (\nabla\vf(\vy_k) - \nabla\vf(\vy_{k-1}))$. By Lemma~\ref{lem:hessian_unbiased}, $\E[\bm\zeta_{H,k} | \mathcal{F}_{k-1}] = \mathbf{0}$. By Assumption~\ref{assump:noise}, $\E[\bm\xi_{g,k} | \mathcal{F}_{k-1}] = \mathbf{0}$. Both are martingale differences, conditionally independent.

\emph{Step 2: Moment bound on $\bm\zeta_{H,k}$.} By Step 1 of the proof of Proposition~\ref{prop:hessian_implies_avg}, $\E\|\nabla^2\tilde\vf(\vx;\zeta)\|_{\mathrm{op}}^r \le \bar L^r$. Hence $\E\|\tilde\mH_k(\vy_k - \vy_{k-1})\|^r \le \bar L^r \gamma^r D_\cX^r$. Since $\bm\zeta_{H,k}$ is a centering of $\tilde\mH_k(\vy_k - \vy_{k-1})$, the standard centering inequality $\E\|\vV - \E\vV\|^r \le 2^{r-1}\E\|\vV\|^r$ gives $\E\|\bm\zeta_{H,k}\|^r \le 2^{r-1}\bar L^r \gamma^r D_\cX^r$.

\emph{Step 3: Conditional von Bahr--Esseen.} Using conditional independence and $(1-\beta)^r \le 1-\beta$:
\[
\E\|\bm\delta_{g,k}\|^r \le (1-\beta)\E\|\bm\delta_{g,k-1}\|^r + 2^{r-1} C_r (1-\beta)\bar L^r \gamma^r D_\cX^r + C_r \beta^r \sigma_g^r.
\]

\emph{Step 4: Iterate.} Unrolling, with $\sum_i (1-\beta)^{k-i} \le 1/\beta$, gives the claim.
\end{proof}

\begin{theorem}[Hessian-corrected acceleration] \label{thm:hessian_rate}
Under Assumptions~\ref{assump:domain}--\ref{assump:noise} and \ref{assump:hessian_var}, with the update \eqref{eq:hessian_grad}, Variant II function tracker, and constant schedule $\gamma = \beta = \rho = K^{-r/(2r-1)}$:
\[
\min_{1 \le k \le K}\E[\hat\Delta_k] = \mathcal{O}\!\left(K^{-(r-1)/(2r-1)}\right).
\]
\end{theorem}

\begin{proof}
By Jensen on Lemma~\ref{lem:hessian_track}, the drift $\gamma\beta^{-1/r}$ and noise $\beta^{(r-1)/r}$ both scale as $K^{-(r-1)/(2r-1)}$. The function tracker inherits this rate via cross-coupling. Telescoping the GLMO progress bound and dividing by $\gamma K$ yields the result.
\end{proof}

\begin{remark}[STORM vs.\ Hessian-corrected]
Both Theorems~\ref{thm:storm_rate} and~\ref{thm:hessian_rate} achieve the rate $\mathcal{O}(K^{-(r-1)/(2r-1)})$ at $r = 2$. STORM reuses $\xi_k$ at two points; the Hessian variant requires an independent Hessian oracle. Proposition~\ref{prop:hessian_implies_avg} shows the latter is strictly stronger, since bounded Hessian moments imply average smoothness. The lower bound of \citet{arjevani2023lower} confirms that the $K^{-1/4}$ rate of Theorem~\ref{thm:nonconvex} is minimax-optimal under Assumption~\ref{assump:inner_smooth} alone; Theorems~\ref{thm:storm_rate} and~\ref{thm:hessian_rate} circumvent this only via the strictly stronger smoothness assumption.
\end{remark}

\section{Additional Examples} \label{sec:examples_appendix}

\begin{example}[Fairness-constrained learning]
Beyond the examples in the main text, Algorithm~\ref{alg:hybrid_sfw} applies to fairness-constrained ERM \citep{agarwal2018reductions, donini2018empirical}, where one minimizes
\[
\min_{\vx \in \cX} f_0(\vx) + \lambda \max_{1 \le i \le n}|f_i(\vx) - f_0(\vx)| + \Xi(\vx).
\]
The outer function $F(\vu, \vx) = u^{(0)} + \lambda\max_i|u^{(i)} - u^{(0)}| + \Xi(\vx)$ depends explicitly on $\vx$. The GLMO reduces to an LP via auxiliary variables for the absolute values.
\end{example}

\begin{example}[Exact penalty for stochastic constraints]
For safe RL or fairness-constrained classification, one can minimize an expected loss subject to expectation constraints via the exact penalty formulation \citep{cotter2019two}:
\[
\min_{\vx \in \cX} f_0(\vx) + \langle \vc, \vx \rangle + \lambda \sum_{i=1}^n \max(0, f_i(\vx) - \tau_i).
\]
The GLMO is again an LP.
\end{example}





\end{document}
